\documentclass[11pt]{article}
\usepackage[T1]{fontenc}
\usepackage{lmodern,microtype}
\usepackage[margin=1in]{geometry}
\usepackage{amsmath,amssymb,amsthm,mathtools,mathrsfs}
\usepackage{booktabs,tabularx,longtable,array,enumitem}
\usepackage[dvipsnames]{xcolor}
\usepackage[colorlinks=true,linkcolor=MidnightBlue,citecolor=MidnightBlue,urlcolor=MidnightBlue]{hyperref}
\usepackage{bookmark}
\usepackage{fancyhdr}
\hypersetup{pdftitle={Riemann--Siegel analysis at a possible positive Newman threshold},
 pdfauthor={Drafted for Edward Baker with OpenAI Codex assistance},
 pdfsubject={Analytical synthesis of RSCT rounds 1--9; 12 November 2026}}
\setlength{\headheight}{14pt}
\pagestyle{fancy}\fancyhf{}
\fancyhead[L]{\small Riemann--Siegel at a possible Newman threshold}
\fancyhead[R]{\small 12 November 2026}
\fancyfoot[C]{\thepage}
\setlist{topsep=5pt,itemsep=3pt}
\setlength{\emergencystretch}{2em}
\numberwithin{equation}{section}
\newtheorem{theorem}{Theorem}[section]
\newtheorem{proposition}[theorem]{Proposition}
\newtheorem{lemma}[theorem]{Lemma}
\newtheorem{corollary}[theorem]{Corollary}
\theoremstyle{definition}
\newtheorem{definition}[theorem]{Definition}
\newtheorem{inputthm}[theorem]{Inherited input}
\theoremstyle{remark}
\newtheorem{remark}[theorem]{Remark}
\newcommand{\R}{\mathbb R}
\newcommand{\C}{\mathbb C}
\newcommand{\N}{\mathbb N}
\newcommand{\dd}{\,\mathrm d}
\newcommand{\ii}{\mathrm i}
\newcommand{\ee}{\mathrm e}
\newcommand{\Log}{\operatorname{Log}}
\newcommand{\Rea}{\operatorname{Re}}
\newcommand{\Ima}{\operatorname{Im}}
\newcommand{\diag}{\operatorname{diag}}
\newcommand{\op}{\mathrm{op}}
\newcommand{\can}{\mathrm{can}}
\newcommand{\vM}{\mathrm{vM}}
\newcommand{\cQ}{\mathcal Q}
\newcommand{\cR}{\mathcal R}
\newcommand{\cB}{\mathcal B}
\newcommand{\cD}{\mathcal D}
\newcommand{\cJ}{\mathcal J}
\newcommand{\norm}[1]{\left\lVert#1\right\rVert}
\providecommand{\RepoRoot}{file:///Users/ebbaker/Documents/shifted-zeta-positivity}
\newcommand{\repo}[2]{\href{\RepoRoot/#1}{#2}}
\newcommand{\rsctnote}[2]{\repo{papers/investigations-5/10_riemann_siegel_critical_time/notes/#1}{#2}}
\title{Riemann--Siegel analysis at a possible\\positive Newman threshold\\[0.5em]
 \large Complete local approximation, signed zero measurements,\\
 and the limits of calibrated phase constraints}
\author{Drafted for Edward Baker}
\date{12 November 2026}
\begin{document}
\maketitle
\begin{abstract}
We develop the analytical results of the Riemann--Siegel critical-time
investigation through its ninth round. A hypothetical positive Newman
threshold has a finite multiple real zero. The complete reflected
Riemann--Siegel formula therefore supplies an arithmetic expression to
compare with necessary inequalities of an all-real zero configuration.
At fixed imaginary heights, a three-height logarithmic estimate forces a
negative noncentral contribution, while a correlated matrix test admits
every fixed multiplicity at sufficiently large centers. At distance
\(\beta/\log(a/(4\pi))\), a complete local approximation and event
subtraction give an absolute signed-measurement error tending to zero.
The remaining difficulty is the sign of the actual correlated arithmetic
expression. We prove, within the inherited approximation framework, that
coefficient disks, exact coefficient circles, and completely
multiplicative unit phases with growing actual-prime calibration all
admit negative values. The latter construction also satisfies the
forced-mirror fourth-derivative inequality with a strict margin.
These are quantitative exclusions of specified estimating classes,
not exclusions of canonical collisions. We finish by identifying the
unpaid canonical-phase, exceptional-center, domain-coverage, and
independent-review tasks.
\end{abstract}

\noindent\textbf{Evidence and authorship.}
This research draft was prepared for Edward Baker with substantial
OpenAI Codex LLM assistance. The recorded session configuration does not
identify the actual backend model or reasoning effort; neither is
inferred. The source status is \emph{proved in-note from named inherited
approximation inputs; same-assistant checking only} for rounds 2--6,
and \emph{proved in-note from named classical and inherited approximation
inputs; same-assistant checking only} for rounds 7--9.
Round 1 supplies elementary identities and a scoped failed acquisition.
The synthesis does not upgrade those statuses or constitute an
independent audit of the inherited field approximation.

\clearpage
{\small\tableofcontents}
\clearpage
\section{The question and the scope of the results}\label{sec:scope}

The investigation asks whether the arithmetic of the complete
Riemann--Siegel representation can rule out the multiple real zero that
would occur at a positive de Bruijn--Newman threshold. There are two
distinct mathematical obligations: obtain a complete, accurately
normalized arithmetic expression, and prove a sign for that expression
that contradicts the necessary geometry of the proposed event.
The first obligation has been met in the local regime developed below.
The second remains open.

The organization follows this distinction. Sections~\ref{sec:geometry}
and~\ref{sec:representation} identify the event and a denominator-free
signed measurement. Section~\ref{sec:fixed} explains why the tested
fixed-height measurements do not resolve a fixed multiplicity.
Section~\ref{sec:micro} derives the microscopic approximation, its
coefficient moments, and the error after imposing the two event equations.
Sections~\ref{sec:phases}--\ref{sec:mirror} test progressively more
restrictive phase information. The final construction keeps all
coefficient magnitudes, complete unit-phase multiplicativity, actual
phases at a growing set of small primes, the actual carrier and correction,
both local event equations, and one necessary fourth-derivative condition.
It still has the sign permitted by a collision.

Throughout the asymptotic sections,
\begin{equation}\label{eq:band}
 L=\log\frac{a}{4\pi},\qquad T=\frac cL,\qquad
 \frac12\le c\le1,\qquad a\longrightarrow\infty.
\end{equation}
Constants are uniform in this interval of \(c\), unless a different
dependence is stated. A fixed disk radius, a fixed \(\beta>0\), or a fixed
derivative order may enter the constant. No bound stated for a fixed
parameter is used with that parameter growing without a separate proof.
The eventual objective concerns every \(0<T\le1/2\), every \(a>0\), and
every multiplicity \(m\ge2\). There is no premise locating a hypothetical
threshold event in \eqref{eq:band}.

The research classification follows the
\repo{result-summaries/README.md}{repository result index}, especially
\repo{result-summaries/11_newman_heat_flow_reductions_and_certificates.md}{cluster 11},
and the \repo{avoid/README.md}{avoidance index}, particularly
\repo{avoid/01_restating_the_target.md}{target restatement},
\repo{avoid/02_shape_only_arguments_and_the_controls.md}{shared control premises},
\repo{avoid/07_scoped_exclusions_read_as_universal.md}{scoped exclusions},
and \repo{avoid/12_term_by_term_matching_and_separate_bounds.md}{separate bounds}.
The endpoint reduction is stated once below. The measurable contributions
are the explicit acquisition errors, moment asymptotics, matrix scale,
and signed counterexamples on precisely defined enlarged classes.

\section{The canonical flow and the finite event}\label{sec:geometry}

\subsection{Normalization and the positive-threshold reduction}

We use the normalization
\begin{align}
 \Phi(u)&=\sum_{n\ge1}
 (2\pi^2n^4\ee^{9u}-3\pi n^2\ee^{5u})
 \ee^{-\pi n^2\ee^{4u}},\label{eq:source}\\
 H_t(z)&=\int_0^\infty \ee^{tu^2}\Phi(u)\cos(zu)\dd u,\qquad
 H_0(z)=\frac18\xi\!\left(\frac{1+\ii z}{2}\right).
 \label{eq:flow}
\end{align}
The rapid source decay permits differentiation on compact sets:
\begin{equation}\label{eq:heat}
 \partial_t H_t=-\partial_{zz}H_t.
\end{equation}
For real \(t\), \(H_t\) is even, real on the real axis, and entire of
order one. Positivity of \(\Phi\) on \(u\ge0\) gives
\(H_t(\ii y)>0\) for real \(y\). The finite de Bruijn--Newman threshold
\(-\infty<\Lambda\le1/2\) satisfies
\begin{equation}\label{eq:threshold}
 \{\text{all zeros of }H_t\text{ are real}\} \quad\Longleftrightarrow\quad
 t\ge\Lambda.
\end{equation}
These classical inputs use the normalization of \cite{polymath}.

\begin{proposition}[Inherited finite-attainment reduction]\label{prop:finite}
If \(\Lambda>0\), then at \(T=\Lambda\) there are a finite \(a>0\)
and an integer \(m\ge2\) such that \(a\) and \(-a\) are zeros of
\(H_T\), each of multiplicity \(m\).
\end{proposition}
\begin{proof}
Take \(t_j\uparrow\Lambda\), with \(t_j\ge\Lambda/2>0\), and choose a
nonreal zero \(z_j\) of \(H_{t_j}\). The positive-time zero-strip theorem
places these zeros in \(|\Ima z|\le1\). The uniform high-zero theorem
\cite[Theorem 1.5(i)]{polymath} makes every zero with
\(|\Rea z|\ge\exp(C/t_j)\) real. Thus, after using the symmetries and
passing to a subsequence, the \(z_j\) lie in one compact set and converge
to a zero \(a\) of \(H_\Lambda\). All zeros at \(\Lambda\) are real, and
\(H_\Lambda(0)>0\), so choose the symmetry with \(a>0\).

If \(a\) were simple, the implicit-function theorem would give a unique
nearby zero for real times near \(\Lambda\). Conjugation would make this
branch real. This contradicts the approaching nonreal zeros. Evenness
supplies the mirror zero with the same multiplicity. This is the
finite-attainment argument already recorded in \cite{collision}.
\end{proof}

For clarity, the endpoint equivalence used by the program is
\begin{equation}\label{eq:endpoint}
 \mathrm{RH}\quad\Longleftrightarrow\quad
 \text{there is no real }a\text{ with }H_t(a)=H_t'(a)=0
 \text{ for any }t>0.
\end{equation}
The reverse implication follows from Proposition~\ref{prop:finite}.
For the forward implication, all times \(t\ge0\) are all-real under RH.
A multiple real zero at \(t>0\) would split nonreally immediately
backwards. Indeed, applying the heat equation to the first nonzero
Taylor term \(b(z-a)^m\) gives
\[
 \ee^{-\tau\partial_z^2}(z-a)^m
 =(2\tau)^{m/2}
 \operatorname{He}_m\!\left(\frac{z-a}{\sqrt{2\tau}}\right).
\]
The simple real zeros of the probabilists' Hermite polynomial give local
branches \(a+\sqrt{2\tau}\,\lambda_j+O(|\tau|)\); for \(\tau<0\),
some are nonreal. This argument, including arbitrary multiplicity, is
inherited from \cite{collision,polymath}. Equation~\eqref{eq:endpoint}
is a target-equivalent reformulation, not an additional estimate.

\subsection{The reflected Poisson measurement}

Fix an all-real slice \(H_T\) and write its even product as
\[
 H_T(z)=H_T(0)\prod_{r>0}\left(1-\frac{z^2}{r^2}\right),
\]
with positive zeros repeated by multiplicity. Pairing the logarithmic
derivative, and then taking its imaginary part, gives an absolutely
convergent sum
\begin{equation}\label{eq:poisson}
 Q(a,y):=-y\Ima\frac{H_T'(a+\ii y)}{H_T(a+\ii y)}
 =\sum_r\frac{y^2}{(a-r)^2+y^2},\qquad y>0.
\end{equation}
Here and below an unqualified sum over \(r\) includes both signs and
multiplicities. The order-one product and its differentiated form are
the inputs used in \cite[Section 3]{polymath}.

Suppose \(a\) has multiplicity \(m\ge2\). Define
\begin{equation}\label{eq:ks}
 k(a,y)=1+\frac{y^2}{4a^2+y^2},\qquad
 S(a,y)=\sum_{r\notin\{a,-a\}}\frac1{(a-r)^2+y^2},
\end{equation}
removing all \(m\) copies of each event root. Then
\begin{equation}\label{eq:eventmass}
 Q(a,y)=m k(a,y)+y^2S(a,y).
\end{equation}
The mirror term is essential even when it is asymptotically small.

\begin{proposition}[Denominator-free necessary sign]\label{prop:sign}
Set, at \(z=a+\ii y\),
\begin{equation}\label{eq:B}
 \cB(a,y)=2k(a,y)|H_T(z)|^2+
             y\Ima\!\left(H_T'(z)\overline{H_T(z)}\right).
\end{equation}
At an all-real canonical event of multiplicity \(m\ge2\),
\begin{equation}\label{eq:Bnegative}
 \cB(a,y)=\bigl((2-m)k(a,y)-y^2S(a,y)\bigr)|H_T(a+\ii y)|^2<0.
\end{equation}
\end{proposition}
\begin{proof}
Multiply \eqref{eq:eventmass} by \(|H_T|^2\) and subtract it from
\(2k|H_T|^2\). All-realness ensures \(H_T(a+\ii y)\ne0\).
The canonical field has further zeros, so \(S>0\).
\end{proof}

At a fixed all-real canonical slice, existence of a \(y>0\) with
\(\cB(a,y)>0\) is equivalent to \(a\) not being a multiple zero.
If \(H_T(a)\ne0\), then \(Q(a,y)\to0\); if \(a\) is simple, then
\(Q(a,y)\to1\). Since \(k\to1\), these cases give positivity for small
\(y\). Multiple zeros are covered by \eqref{eq:Bnegative}.
This observation identifies the strength of an unrestricted sign demand;
the research task is to obtain a concrete arithmetic estimate on a
declared domain.

\subsection{What an accurate arithmetic value would have to pay}

Let \(F=H_T(z)\), and suppose \(f,g\) approximate its complete value
and physical derivative:
\[
 |F-f|\le\epsilon_0,\qquad |F'-g|\le\epsilon_1.
\]
Expanding each product yields
\begin{equation}\label{eq:Berror}
\begin{split}
 |\cB-(2k|f|^2+y\Ima(g\bar f))|
 \le{}&2k(2|f|\epsilon_0+\epsilon_0^2)\\
 &+y\bigl(|g|\epsilon_0+|f|\epsilon_1+\epsilon_0\epsilon_1\bigr).
\end{split}
\end{equation}
Thus a positive computed expression must exceed this entire error.
No denominator floor is required. If a quotient is used instead,
\(|f|>\epsilon_0\) is necessary for the bound
\begin{equation}\label{eq:quotient}
 \left|Q+y\Ima\frac gf\right|
 \le y\,\frac{\epsilon_1+|g/f|\epsilon_0}{|f|-\epsilon_0}.
\end{equation}

The cancellation at a double zero explains the sensitivity of this test.
Write \(H_T(a+w)=b_mw^m+O(w^{m+1})\), \(b_m\ne0\), and put
\(S_0=\sum_{r\ne\pm a}(a-r)^{-2}>0\).
Isolation of the event and convergence of the squared reciprocal tail
give \(S(a,y)\to S_0\). Hence
\begin{equation}\label{eq:small-y}
 \cB(a,y)=
 \begin{cases}
 -|b_2|^2 S_0 y^6(1+o(1)),&m=2,\\
 (2-m)|b_m|^2 y^{2m}(1+o(1)),&m>2.
 \end{cases}
\end{equation}
For \(m=2\), the separate terms in \eqref{eq:B} have size \(y^4\).
In the separate absolute envelope \eqref{eq:Berror},
\(\epsilon_0=o(y^4)\) and \(\epsilon_1=o(y^3)\) suffice for an
\(o(y^6)\) error. These are fixed-event sufficient conditions for that
envelope; they are not lower bounds on actual approximation errors.

\section{The complete reflected representation}\label{sec:representation}

\subsection{An exact identity before asymptotics}

Let
\[
 s_\pm(z)=\frac{1\pm\ii z}{2},\qquad
 M(s)=\frac{s(s-1)}{16}\pi^{-s/2}\Gamma(s/2),\qquad
 U^*(s)=\overline{U(\bar s)}.
\]
The exact Riemann--Siegel contour remainder is
\begin{equation}\label{eq:contour}
 R_{0,N}(s)=M(s)\int_{N\swarrow N+1}
       \frac{w^{-s}\ee^{\ii\pi w^2}}{\ee^{\pi\ii w}-\ee^{-\pi\ii w}}\dd w.
\end{equation}
The contour crosses \((N,N+1)\) and has the orientation specified in
\cite[Section 4, equations (36)--(38)]{polymath}. For \(T>0\) and
nonreal \(s\), put
\begin{align}
 r_{T,n}(s)&=\frac1{\sqrt\pi}\int_\R
  \ee^{-v^2}M(s+\sqrt T\,v)n^{-s-\sqrt T\,v}\dd v,\\
 R_{T,N}(s)&=\frac1{\sqrt\pi}\int_\R
  \ee^{-v^2}R_{0,N}(s+\sqrt T\,v)\dd v,\qquad
 U_{T,N}(s)=\sum_{n\le N}r_{T,n}(s)+R_{T,N}(s).
\end{align}
For every fixed integer \(N\ge0\) and \(\Rea z\ne0\),
\begin{equation}\label{eq:exactRS}
 H_T(z)=U_{T,N}(s_+(z))+U_{T,N}^*(s_-(z)).
\end{equation}
To see the heat step, insert the Gaussian identity
\[
 \ee^{Tu^2}\cos(zu)=\frac1{\sqrt\pi}\int_\R
       \ee^{-v^2}\cos((z-2\ii\sqrt T\,v)u)\dd v
\]
in \eqref{eq:flow}, and then insert the exact two-contour formula for
\(H_0\). The Gaussian dominates the sub-Gaussian growth of each
complete arm on the permitted shifted lines. The same local domination
justifies differentiation:
\begin{equation}\label{eq:exactRSprime}
 H_T'(z)=\frac\ii2U_{T,N}'(s_+(z))
              -\frac\ii2(U_{T,N}^*)'(s_-(z)).
\end{equation}
Both the gamma factor and the full contour remainder are differentiated.
Off the real axis, \eqref{eq:exactRS} is not twice a real part.

On the real axis, let \(\mathscr A(x)=U_{T,N}(s_+(x))\).
Then \(H_T=2\Rea\mathscr A\). If \(\mathscr A=\rho\ee^{\ii\psi}\ne0\)
at a zero, differentiation gives
\[
 H_T'=-2\rho\psi'\sin\psi.
\]
A multiple zero therefore requires \(\psi'=0\), or belongs to the
separate alternative \(\mathscr A=0\). No strict phase rotation or
nonzero arm floor follows from the reflected representation alone.
Indeed \(\mathscr A\mapsto\mathscr A+\ii q\), with \(q\) a real-entire
heat solution, preserves the reduced data \(H=\mathscr A+\mathscr A^*\).
This is an information test, not permission to alter the canonical arm.
Section~\ref{sec:controls} gives an exact threshold control.

\subsection{The named approximation input and its domain}

The asymptotic work uses the internal complex-strip theorem HDA CS3/CS6
and its application GNP-RS E01 \cite{hda,gnpe01}.
We state exactly the portion consumed here. It is not re-proved as an
independent provider theorem in this manuscript.

On a fixed chart about a large positive real center \(a\), choose the
branches continued from \(a\), and hold
\(N=\lfloor\sqrt{a/(4\pi)}\rfloor\) fixed. Define
\begin{align}
 L(z)&=\Log(z/(4\pi)),\\
 A(T,z)&=\frac{\pi^{1/4}}{16}z^{7/4}
   \exp\left[-\frac{\pi z}{8}
       +\frac T{16}\left(L(z)^2-\frac{\pi^2}{4}\right)\right],
       \label{eq:carrier}\\
 \Theta(T,z)&=\frac z4(L(z)-1)+\frac{\pi T L(z)}{16}
                         +\frac{7\pi}{8},\qquad
 \nu(z)=\frac z2+\frac{\pi T}{8},\\
 C_\pm(T,z)&=\frac{A(T,z)}2\ee^{\pm\ii\Theta(T,z)},\\
 w_n(T,z)&=n^{-1/2}
  \exp\left\{\frac T4\bigl((\log n)^2-L(z)\log n\bigr)\right\}.
\end{align}
The retained correction is
\begin{align}
 \rho(z)&=\sqrt{z/(4\pi)},\qquad
 G(z)=\bigl[M(s_+(z))M(s_-(z))\bigr]^{1/2},\\
 \mathcal C(p)&=\frac{\cos(2\pi(p^2-p)-\pi/8)}{\cos(2\pi p)},\qquad
 \mathcal R_N(z)=(-1)^N G(z)\rho(z)^{-1/2}\mathcal C(\rho(z)-N).
 \label{eq:Rcorrection}
\end{align}
The branch of \(G\) is positive on the real segment. Apparent poles of
\(\mathcal C\) are removable. This is the exact-gamma correction in the
approximant; it must not be confused with the full exact contour
remainder \(R_{T,N}\) in \eqref{eq:exactRS}.
Put
\begin{equation}\label{eq:FN}
 F_N(T,z)=A(T,z)\sum_{n\le N}w_n(T,z)
       \cos\bigl(\Theta(T,z)-\nu(z)\log n\bigr)
       +\ee^{-\pi^2T/64}\mathcal R_N(z).
\end{equation}

\begin{inputthm}[Buffered complete field estimate]\label{input:provider}
For each declared fixed target rectangle, there are finite constants and
an onset such that, for \(0\le T\log(a/(4\pi))\le8\),
\begin{equation}\label{eq:provider}
 |H_T(z)-F_N(T,z)|\le C a^{-3/4}(|C_+(T,z)|+|C_-(T,z)|).
\end{equation}
The difference is holomorphic on the chart. The input is uniform in
fractional-cutoff position. It includes the two complete populations,
the retained correction, the enlarged source charts, and the entire
Gaussian exterior.
\end{inputthm}

For the present application one may take target horizontal radius \(1\)
and vertical height \(9\), retaining the larger source buffers of
\cite{gnpe01}. Its tail constants \(c_9,P_9\) are fixed before choosing
\[
 V_9\ge\max\{16,17,16c_9,8\sqrt{P_9+6}\}.
\]
The required initial charts have radius \(V_9+3\), height \(10\), and
radius \(2\), height \(10+2V_9\), with their distinct onsets.
This records the dependence that permits Cauchy estimates later.
No unbounded Gaussian integral is estimated from a bound valid only on
a narrow strip.

Polymath's effective value Theorem 1.3 has its own printed domain
\(x\ge200,\ 0\le y\le1,\ 0<T\le1/2\). We do not enlarge that domain
or differentiate that theorem to obtain our buffered estimates.
The broader complex-strip statement is the named internal input above,
at its existing evidence level.

\section{Fixed-height measurements and the noncentral zeros}\label{sec:fixed}

\subsection{A signed three-height arithmetic estimate}

At the parameters \eqref{eq:band}, consider
\begin{equation}\label{eq:threeheight}
 W_T(a)=2\log|H_T(a+6\ii)|
              -\log|H_T(a+4\ii)|-\log|H_T(a+8\ii)|.
\end{equation}
Normalize by \(V=H_T/C_-\), and write
\(\nu=a/2+\pi T/8\).

\begin{lemma}[Complete outer comparison]\label{lem:outer}
For \(y=4,6,8\), the inherited field input implies, after a finite
uniform onset,
\begin{equation}\label{eq:outer}
 \left|V(a+\ii y)-
 \zeta\left(\frac{1+y}{2}+\frac c4-\ii\nu\right)\right|
 \le\frac{C_F^*}{L}.
\end{equation}
The constant \(C_F^*\) is finite and independent of \(a,c\).
\end{lemma}
\begin{proof}
Divide \eqref{eq:FN} by \(C_-\). The dominant exponential arm is a
Dirichlet sum whose argument is
\((1-\ii z)/2+T L(z)/4-\ii\pi T/8\), with the additional quadratic
weight \(\exp(T(\log n)^2/4)\).
The calculation in GNP-RS E01 \cite{gnpe01}, extended over the finite
set \(\{4,6,8\}\), has the following complete ledger:
\[
\begin{array}{ll}
\text{reflected population}&O(a^{(1-y)/4}),\\
\text{retained correction}&O(a^{-(y+1)/4}),\\
\text{complete normalized field error}&O(a^{-3/4}),\\
\text{quadratic-weight removal}&\le T D_2/4,\\
\text{ordinary Dirichlet tail}&\le(2/3)N^{-3/2},\\
\text{argument replacement}&\le D_1T|L(a+\ii y)-L|/4,
\end{array}
\]
where \(D_j=\sum_{n\ge2}n^{-5/2}(\log n)^j<\infty\).
For the weight removal, \(T\log n/4\le c/8\), and the elementary
bound \(|\ee^v-1|\le v\ee^v\) leaves a summable majorant.
Every other term is \(O(1/L)\) or smaller, uniformly at the three
heights. This yields \eqref{eq:outer}; the population is summed before
its error is bounded.
\end{proof}

Since
\[
 |\zeta(s)-1|\le 2^{-5/2}+\int_2^\infty x^{-5/2}\dd x
               =\frac7{12\sqrt2}<\frac12
 \quad(\Rea s\ge5/2),
\]
the gate \(L>4C_F^*\) supplies \(|V|>1/4\). Integration of the
logarithmic derivative along the straight segment from \(\zeta(s)\) to
\(V\) gives
\begin{equation}\label{eq:outerlog}
 \bigl|\log|V|-\log|\zeta(s)|\bigr|\le4C_F^*/L.
\end{equation}
This is an off-axis floor at these observations only.

Let \(g(z)=\log C_-(T,z)\), with the local analytic logarithm. Then
\begin{align}
 g'(z)&=\frac7{4z}-\frac\pi8+\frac{TL(z)}{8z}
              -\frac{\ii L(z)}4-\frac{\ii\pi T}{16z},\label{eq:gprime}\\
 g''(z)&=-\frac7{4z^2}+\frac{T(1-L(z))}{8z^2}
              -\frac\ii{4z}+\frac{\ii\pi T}{16z^2}.
\end{align}
For \(u(y)=\log|C_-(T,a+\ii y)|\), \(u''=-\Rea g''\).
For large \(a\), \(T|L(a+\ii y)|\le2\), and
\(|\Rea(-\ii/(a+\ii y))|=y/(a^2+y^2)\). Consequently
\[
 |u''(y)|\le5/a^2\quad(4\le y\le8),\qquad
 |D_C|\le20/a^2,\quad D_C=2u(6)-u(4)-u(8).
\]

\begin{proposition}[Three-height bound]\label{prop:W}
Let \(q_n=\Lambda_{\vM}(n)/\log n\), where \(\Lambda_{\vM}\) is the
von Mangoldt function. Then
\begin{equation}\label{eq:Wseries}
 W_T(a)=-\sum_{n\ge2}q_n n^{-5/2-c/4}(1-n^{-1})^2
                    \cos(\nu\log n)+D_C+e,\qquad
 |e|\le16C_F^*/L.
\end{equation}
In particular,
\begin{equation}\label{eq:Wbound}
 |W_T(a)|<\frac16+\frac{16C_F^*}{L}+\frac{20}{a^2}.
\end{equation}
Under the eventual gates
\begin{equation}\label{eq:Wgates}
 16C_F^*/L+20/a^2\le1/72,\qquad a^2\ge288,
\end{equation}
one has \(|W_T(a)|<13/72\).
\end{proposition}
\begin{proof}
The absolutely convergent Euler logarithm of ordinary zeta in
\(\Rea s>1\) \cite{dlmfzeta} is
\(\log\zeta(s)=\sum_{n\ge2}q_n n^{-s}\).
Combine its three values before taking absolute values. Their
coefficient is \(2n^{-1}-1-n^{-2}=-(1-n^{-1})^2\), with the same
phase \(\nu\log n\) throughout. The four total units of logarithmic
error in \eqref{eq:outerlog} give \eqref{eq:Wseries}.
The function \(f(x)=x^{-5/2}(1-x^{-1})^2\) decreases for \(x\ge2\);
since \(0\le q_n\le1\),
\[
 \sum_{n\ge2}q_n n^{-5/2-c/4}(1-n^{-1})^2
 \le f(2)+\int_2^\infty f(x)\dd x
 =\frac{389}{1680\sqrt2}<\frac16.
\]
This proves the claimed bound. A truncation at an integer \(M\ge2\)
would have tail at most
\(M^{-3/2-c/4}/(3/2+c/4)\); no truncation or sampled value is needed.
\end{proof}

\subsection{Why the scalar bound requires a negative background}

At an all-real slice the same three-height combination of the product is
\begin{equation}\label{eq:K}
 W_T(a)=\sum_r K(a-r),\qquad
 K(d)=\frac12\log\frac{(d^2+36)^2}{(d^2+16)(d^2+64)}.
\end{equation}
The sum converges absolutely because \(K(d)=-4d^{-2}+O(d^{-4})\).
At an event of multiplicity \(m\ge2\),
\begin{equation}
 W_T(a)=m[K(0)+K(2a)]+\mathcal R_{\mathrm{bg}},\qquad
 \mathcal R_{\mathrm{bg}}=\sum_{r\ne\pm a}K(a-r).
 \label{eq:Wbackground}
\end{equation}
Now \(K(0)=\log(9/8)>1/9\). For \(d^2\ge64\), direct comparison of
the numerator and denominator gives \(|K(d)|\le8/d^2\).
Thus \eqref{eq:Wgates} implies
\[
 m[K(0)+K(2a)]>2(1/9-1/144)=15/72.
\]
Combining with Proposition~\ref{prop:W} gives the quantitative consequence
\begin{equation}\label{eq:bgnegative}
 \boxed{\mathcal R_{\mathrm{bg}}<-1/36.}
\end{equation}
It is a necessary signed background contribution under the event,
not a contradiction: \(K\) changes sign.

\begin{proposition}[Cancellation forces a sign-changing log kernel]
\label{prop:kernel}
Let \(b_j>0\) and \(\alpha_j\in\R\) be a finite collection, and put
\[
 K_\alpha(d)=\frac12\sum_j\alpha_j\log(d^2+b_j^2).
\]
If \(\sum_j\alpha_j=0\), this kernel is integrable and
\(\int_\R K_\alpha(d)\dd d=\pi\sum_j\alpha_j b_j\).
If also \(\sum_j\alpha_j b_j=0\) and \(K_\alpha(0)>0\), then
\(K_\alpha\) is negative somewhere.
\end{proposition}
\begin{proof}
The first cancellation removes the \(\log d^2\) term, leaving
\(O(d^{-2})\). For positive \(b,b_0\), differentiating in \(b\) gives
\[
 \int_\R\log\frac{d^2+b^2}{d^2+b_0^2}\dd d=2\pi(b-b_0).
\]
Form the linear combination to obtain the integral. A continuous,
nonnegative, integrable function that is positive at zero cannot have
zero integral.
\end{proof}

For \eqref{eq:K},
\[
 K(d)>0\ (|d|<\sqrt{34}),\qquad
 K(d)<0\ (|d|>\sqrt{34}),\qquad \int_\R K(d)\dd d=0.
\]
The proposition concerns finite fixed-height logarithmic combinations
that cancel the linear carrier. It does not rule out a signed arithmetic
estimate of the background, explicit deflation, or a different observable.
The genuine positive-source threshold control in
Section~\ref{sec:controls} satisfies the same scalar cap.

\subsection{A correlated matrix test}

Write \(h=-H_T'/H_T\), and let \(z_j=a+\ii b_j\), with
\((b_1,b_2,b_3)=(4,6,8)\). On an all-real slice,
\begin{equation}\label{eq:Pick}
 M_{ij}=\frac{h(z_i)-\overline{h(z_j)}}{z_i-\bar z_j}
       =\sum_r\frac1{(z_i-r)(\bar z_j-r)}.
\end{equation}
This is a Gram matrix. Set
\[
 u_j=\frac1{\ii b_j},\qquad v_j=\frac1{2a+\ii b_j},\qquad
 E=uu^*+vv^*.
\]
An event of multiplicity \(m\) necessarily satisfies
\begin{equation}\label{eq:Pickevent}
 M-mE\succeq0.
\end{equation}
This retains correlations among all three observations.

There is a direct link to the previous scalar background. Deflate the
event in the logarithmic derivative:
\[
 h_{\mathrm{bg}}(z)=h(z)+\frac m{z-a}+\frac m{z+a},\qquad
 p_{\mathrm{bg}}(y)=\Ima h_{\mathrm{bg}}(a+\ii y)\ge0.
\]
Integration of \(\partial_y\log|H_T(a+\ii y)|=\Ima h(a+\ii y)\)
gives
\[
 \mathcal R_{\mathrm{bg}}
   =\int_4^6 p_{\mathrm{bg}}(y)\dd y-\int_6^8p_{\mathrm{bg}}(y)\dd y.
\]
Positivity alone does not order these two integrals.

To evaluate \(M\) arithmetically, use disks
\(|z-z_j|\le1/2\) in the fixed provider rectangle. Set
\[
 Z(z)=\zeta\left(\frac{1-\ii z}{2}+\frac c4-\frac{\ii\pi T}{8}\right),
 \qquad \Delta=V-Z.
\]
The complete outer ledger above, now on \(7/2\le\Ima z\le17/2\),
has reflected cost \(O(a^{-5/8})\), correction \(O(a^{-9/8})\),
field error \(O(a^{-3/4})\), summable weight-removal cost \(O(1/L)\),
tail \(O(N^{-5/4})\), and argument cost \(O(1/(aL))\). Hence
\begin{equation}\label{eq:deltaderivative}
 \norm{\Delta}_{|z-z_j|\le1/2}\le C_\Delta/L,\qquad
 |\Delta'(z_j)|\le2C_\Delta/L.
\end{equation}
The second bound is Cauchy's estimate for a holomorphic difference.
Since \(\Rea s\ge9/4\),
\[
 |Z-1|\le \frac{13}{5\,2^{9/4}}<\frac34,\qquad
 |Z'|\le\frac12\sum_{n\ge2}n^{-9/4}\log n=:B_Z.
\]
The gate \(C_\Delta/L\le1/8\) gives \(|V|>1/8\), and therefore
\begin{equation}\label{eq:logdererror}
 \left|\frac{V'}V-\frac{Z'}Z\right|
 \le\frac{C_h}{L},\qquad C_h=(16+32B_Z)C_\Delta.
\end{equation}
All denominators are paid before the matrix is formed.

Put \(\sigma_j=(1+b_j)/2+c/4\) and
\(P_j=\sum_{n\ge2}\Lambda_{\vM}(n)n^{-\sigma_j}\ee^{\ii\nu\log n}\).
Then
\begin{align}
 h(z_j)&=-g'(z_j)-\frac\ii2 P_j+O(1/L),\\
 M_{ij}&=\frac{-g'(z_i)+\overline{g'(z_j)}}{\ii(b_i+b_j)}
       -\frac{P_i+\overline{P_j}}{2(b_i+b_j)}
       +O\!\left(\frac1{L(b_i+b_j)}\right).
 \label{eq:signedM}
\end{align}
The bounded prime-power sums and \eqref{eq:gprime} imply
\begin{equation}\label{eq:Mbulk}
 M=\frac L2 C+O_{\op}(1),\qquad C_{ij}=\frac1{b_i+b_j}.
\end{equation}
This is a statement about the exact canonical samples at every large
center in the band, whether or not it is an event. The leading matrix
is positive definite because
\[
 \alpha^*C\alpha=\int_0^\infty
             \left|\sum_j\alpha_j\ee^{-b_jt}\right|^2\dd t>0
 \quad(\alpha\ne0).
\]

\begin{proposition}[The fixed-height matrix admits fixed multiplicities]
\label{prop:matrixcap}
For large \(a\), the largest formal nonnegative event weight allowed by
\eqref{eq:Pickevent} is
\begin{equation}\label{eq:matrixcap}
 \mathfrak m(a,c)=
 \frac1{\lambda_{\max}(M^{-1/2}EM^{-1/2})}
 =\frac6{13}L+O(1).
\end{equation}
For every fixed integer \(m\ge1\), the exact matrix \(M-mE\) is positive
definite at every sufficiently large center in the band.
\end{proposition}
\begin{proof}
Let \(r_j=1/b_j\). For any distinct positive heights, the Cauchy matrix
satisfies
\begin{equation}\label{eq:Cauchyinverse}
 r^tC^{-1}r=2\sum_j b_j^{-1}.
\end{equation}
To prove it, set \(x=C^{-1}r\) and \(f(z)=\sum_jx_j/(z+b_j)\).
Since \(f(b_i)=1/b_i\),
\[
 1-zf(z)=(-1)^q\prod_{j=1}^q\frac{z-b_j}{z+b_j}.
\]
Differentiate at zero. The left derivative is \(-f(0)\); the right
derivative is \(-2\sum b_j^{-1}\), which proves the identity.

Equation~\eqref{eq:Mbulk} gives
\(M^{-1}=2C^{-1}/L+O_{\op}(L^{-2})\).
The \(u\) term has its single nonzero eigenvalue
\(u^*M^{-1}u=4\sum b_j^{-1}/L+O(L^{-2})\).
The \(v\) term has norm \(O(a^{-2}/L)\). Since
\(\sum_{j=1}^3b_j^{-1}=13/24\), inversion proves
\eqref{eq:matrixcap}. Finally
\[
 \lambda_{\min}(M-mE)
 \ge \tfrac L2\lambda_{\min}(C)-O(1)-m\norm{E}>0
\]
for fixed \(m\) and large \(L\).
\end{proof}

This does not construct a zero configuration realizing the formal event
weight. It shows that this necessary matrix condition cannot reject a
fixed multiplicity in the tested regime. Its scale has a simple
interpretation:
\[
 \int_\R\frac{\dd r}{(a+\ii b_i-r)(a-\ii b_j-r)}
       =\frac{2\pi}{b_i+b_j}.
\]
A continuous density \(L/(4\pi)\) produces exactly the bulk \(LC/2\).
The fixed observations retain a growing background while the proposed
event has fixed weight. Substituting \(b_j=\beta_j/L\) into
\eqref{eq:Mbulk} is not justified by its fixed-geometry error estimate;
the microscopic problem requires a new acquisition argument.

\section{Microscopic acquisition with the complete population}\label{sec:micro}

\subsection{Freezing one analytic chart}

Use \eqref{eq:band} and the scaled coordinate
\[
 z=a+w/L,\qquad N=\lfloor\ee^{L/2}\rfloor.
\]
Time, cutoff, and all center data are held fixed in every \(w\) or
\(z\) derivative. Define
\begin{align}
 d_n&=n^{-1/2}\exp\left\{\frac{c}{4L}
                       [(\log n)^2-L\log n]\right\},\\
 \lambda_n&=\frac14-\frac{\log n}{2L}\in[0,1/4],\qquad
 \phi_n=\Theta(T,a)-\nu(a)\log n,\label{eq:frozendata}\\
 P_a(w)&=\sum_{n\le N}d_n\cos(\phi_n+\lambda_nw),\\
 K_a(w)&=\frac{\ee^{-\pi^2T/64}\mathcal R_N(a+w/L)}
                         {A(T,a+w/L)},\qquad
 q_a=P_a+K_a,\qquad f_a(w)=\frac{H_T(a+w/L)}{A(T,a+w/L)}.
 \label{eq:fq}
\end{align}
The nonvanishing carrier \(A\) is a normalization, not a lower bound for
the field. Both reflected arms are in the cosine, and \(K_a\) is retained
as a whole analytic function.

The positive coefficient moments supply both the freezing cost and the
later error payments.
\begin{lemma}[Fixed-order coefficient moments]\label{lem:moments}
Let \(\delta=1/4-c/16\in[3/16,7/32]\). For every fixed integer \(k\ge0\),
\begin{equation}\label{eq:moments}
 M_k:=\sum_{n\le N}d_n\lambda_n^k
 =\frac{2k!}{L^k}\ee^{\delta L}\bigl(1+O_k(L^{-1})\bigr).
\end{equation}
For \(k=0\), the convention is \(\lambda_n^0=1\), including a possible
\(\lambda_N=0\).
\end{lemma}
\begin{proof}
Extend the summand to \(1\le x\le\ee^{L/2}\). It is nonnegative and
decreasing, because
\[
 \frac{\dd\log d(x)}{\dd\log x}
 =-\frac12-\frac c4+\frac{c\log x}{2L}\in[-3/4,-1/2]
\]
and \(\lambda(x)\) decreases. Integral comparison, including the final
fractional interval, has discrepancy \(O_k(1)\). The change of variables
\(v=L/2-\log x\) gives the exact integral
\[
 \frac{\ee^{\delta L}}{(2L)^k}
 \int_0^{L/2}v^k\exp\left(-\frac v2+\frac{cv^2}{4L}\right)\dd v.
\]
On this interval the exponential is at most \(\ee^{-3v/8}\), and its
difference from \(\ee^{-v/2}\) is at most
\(v^2\ee^{-3v/8}/(4L)\). The integral is therefore
\[
 \int_0^\infty v^k\ee^{-v/2}\dd v+O_k(L^{-1})
 =2^{k+1}k!+O_k(L^{-1}).
\]
The exponential scale absorbs the \(O_k(1)\) discretization error,
uniformly since \(\delta\ge3/16\).
\end{proof}

\begin{theorem}[Complete local approximation]\label{thm:local}
For every fixed \(R>1\), and for sufficiently large \(a\), the inherited
field input gives
\begin{equation}\label{eq:localerror}
 \sup_{|w|\le R}|\partial_w^j(f_a-q_a)|
       \le C_Ra^{-3/4},\qquad 0\le j\le4.
\end{equation}
Moreover, with \(\chi=a^{-1/4}\ee^{-cL/16}\),
\begin{equation}\label{eq:Ksmall}
 \sup_{|w|\le R}|\partial_w^jK_a(w)|\le C_R\chi,
                    \qquad 0\le j\le4.
\end{equation}
The constants and onsets are finite but not numerically instantiated.
\end{theorem}
\begin{proof}
Work first on \(|w|\le2R\), increasing that fixed radius if necessary
for a Cauchy margin. Eventually this lies inside the declared fixed
provider target. Since
\[
 \Ima\Theta(T,a+w/L)=\Ima w/4+O_R((aL^2)^{-1}),
\]
the ratio \((|C_+|+|C_-|)/|A|\) is bounded by a constant depending on
\(R\). Thus \eqref{eq:provider}, after division by \(A\), is
\(O_R(a^{-3/4})\).

Uniformly for \(n\le N\),
\begin{align*}
 L(a+w/L)-L&=O_R((aL)^{-1}),\\
 w_n(T,a+w/L)/d_n&=1+O_R((aL)^{-1}),\\
 \Theta(T,a+w/L)-\nu(a+w/L)\log n
    &=\phi_n+\lambda_nw+O_R((aL^2)^{-1}).
\end{align*}
The last estimate follows from
\(\Theta_z=L(z)/4+\pi T/(16z)\) and \(\Theta_{zz}=O(1/a)\).
Because \(|\lambda_n w|\le R/2\) on the larger disk, summing the
holomorphic freezing error costs at most
\[
 O_R\!\left(\frac{M_0}{aL}\right)
 =O_R(a^{-25/32}/L)=O_R(a^{-3/4}).
\]
The heat correction in the phase and the quadratic weights have been
frozen at their actual values, not discarded.

Apply Cauchy's estimate to this holomorphic difference on the larger
disk. In particular, if its supremum there is \(C_{2R}a^{-3/4}\),
then at any point with a radius-\(R\) Cauchy circle available the
order-\(j\) bound is \(j!R^{-j}C_{2R}a^{-3/4}\).
This proves \eqref{eq:localerror} for the stated finite set of orders.

Finally the correction input gives
\(\mathcal R_N(z)/A(0,z)=O_R(a^{-1/4})\), and exactly
\[
 \ee^{-\pi^2T/64}\frac{A(0,z)}{A(T,z)}
                 =\ee^{-T L(z)^2/16}.
\]
Its modulus on the disk is \(O_R(\ee^{-cL/16})\).
Another Cauchy estimate proves \eqref{eq:Ksmall}. No moving floor or
real-line big-\(O\) term has been differentiated.
\end{proof}

The inherited terminal population scale and earlier phase-retaining
first-derivative approximation belong to \cite{collision,s04}.
The RSCT applications are the fixed-\(k\) constants in
\eqref{eq:moments}, the event-subtracted acquisition, the particular
signed error below, and the fourth-derivative use in
Section~\ref{sec:mirror}.

\subsection{Subtract the necessary event equations first}

For an analytic function \(g\), write
\[
 g^\circ(w)=g(w)-g(0)-wg'(0).
\]
At any actual event of multiplicity \(m\ge2\), the nonzero carrier gives
\(f_a(0)=f_a'(0)=0\). Set \(S_a=q_a^\circ\). If \(e=f_a-q_a\),
then \(f_a-S_a=e-e(0)-we'(0)\), so Taylor's integral remainder yields
\begin{equation}\label{eq:eventerror}
 |f_a(w)-S_a(w)|\le C_Ra^{-3/4}|w|^2,\qquad
 |f_a'(w)-S_a'(w)|\le C_Ra^{-3/4}|w|.
\end{equation}
It also gives \(q_a(0),q_a'(0)=O(a^{-3/4})\).
These are just two necessary equations, valid for every \(m\ge2\);
the remaining \(m-2\) equations have not been imposed.
Explicitly,
\begin{equation}\label{eq:subtracted}
\begin{split}
 S_a(w)={}&\sum_{n\le N}d_n
 \{\cos(\phi_n+\lambda_nw)-\cos\phi_n
                       +\lambda_nw\sin\phi_n\}\\
 &+K_a(w)-K_a(0)-wK_a'(0).
\end{split}
\end{equation}

\begin{proposition}[Size of the separate absolute allowances]\label{prop:allowance}
For fixed \(\beta>0\), define
\[
 U_\beta=\sum_{n\le N}d_n(\cosh(\beta\lambda_n)-1),\qquad
 V_\beta=\sum_{n\le N}d_n\lambda_n\sinh(\beta\lambda_n).
\]
Then
\begin{equation}\label{eq:UV}
 U_\beta=\frac{2\beta^2}{L^2}\ee^{\delta L}(1+O_\beta(L^{-1})),
 \qquad
 V_\beta=\frac{4\beta}{L^2}\ee^{\delta L}(1+O_\beta(L^{-1})).
\end{equation}
These bound the value and derivative in \eqref{eq:subtracted},
apart from the \(O_\beta(\chi)\) correction.
\end{proposition}
\begin{proof}
At \(w=\ii\beta\), a subtracted cosine equals
\[
 \cos\phi\,(\cosh x-1)+\ii\sin\phi\,(x-\sinh x),
 \qquad x=\beta\lambda_n.
\]
Because \(\cosh x-1\ge\sinh x-x\ge0\), its maximum modulus over a free
real phase is \(\cosh x-1\). The derivative has maximal modulus
\(\lambda_n\sinh x\). Taylor's formula, with
\(0\le x\le\beta/4\), gives
\[
 U_\beta=\frac{\beta^2}{2}M_2+O_\beta(M_4),\qquad
 V_\beta=\beta M_2+O_\beta(M_4).
\]
Insert Lemma~\ref{lem:moments}.
\end{proof}

The allowances diverge uniformly in \(c\). Event subtraction saves
\(L^2\), but a separate absolute estimate still pays the terminal
population. This is a lower bound on the size of a specified upper
allowance. It is not a lower bound on \(|S_a|\), on its signed quadratic
form, or on the actual error. In particular, the independent maximizing
phases need not satisfy the event equations or the canonical phase law.

\subsection{A small acquisition error for the signed expression}

Fix \(y=\beta/L\), and define
\begin{align}
 k_a&=1+\frac{(\beta/L)^2}{4a^2+(\beta/L)^2},\\
 \mathfrak d_a&=2k_a+\frac\beta L
              \Ima\frac{A_z(T,a+\ii\beta/L)}{A(T,a+\ii\beta/L)},\\
 \cD_a(g)&=\mathfrak d_a|g(\ii\beta)|^2
                   +\beta\Ima(g'(\ii\beta)\overline{g(\ii\beta)}).
 \label{eq:D}
\end{align}
The carrier logarithmic derivative is
\begin{equation}\label{eq:bcarrier}
 b(z):=\frac{A_z}{A}=\frac7{4z}-\frac\pi8+\frac{TL(z)}{8z},
 \qquad \mathfrak d_a=2+O_\beta(a^{-2}L^{-2}).
\end{equation}
Since \(H_T'=A_zf_a+ALf_a'\), the normalization is exact:
\begin{equation}\label{eq:physicalD}
 \cD_a(f_a)=\frac{\cB(a,\beta/L)}{|A(T,a+\ii\beta/L)|^2}.
\end{equation}
At an all-real event,
\begin{equation}\label{eq:physicalDnegative}
 \cD_a(f_a)=
 \left[(2-m)k_a-\frac{\beta^2}{L^2}
  \sum_{r\ne\pm a}\frac1{(a-r)^2+\beta^2/L^2}\right]|f_a(\ii\beta)|^2<0.
\end{equation}

\begin{proposition}[Complete signed-error payment]\label{prop:Derror}
Under an actual event, \(\cD_a(f_a)\) differs from
\(\cD_a(S_a)\) by at most
\begin{equation}\label{eq:eta}
 \eta_a=O_\beta\left(
       a^{-3/4}\ee^{\delta L}/L^2+a^{-3/2}\right)
       =o(L^{-2}).
\end{equation}
\end{proposition}
\begin{proof}
If the value and derivative errors at \(\ii\beta\) are
\(\epsilon_0,\epsilon_1\), expansion gives
\begin{align*}
 |\cD_a(f_a)-\cD_a(S_a)|
 \le{}&|\mathfrak d_a|(2|S_a|\epsilon_0+\epsilon_0^2)\\
 &+\beta(|S_a'|\epsilon_0+|S_a|\epsilon_1+\epsilon_0\epsilon_1).
\end{align*}
Use \eqref{eq:eventerror}, Proposition~\ref{prop:allowance}, and
\eqref{eq:Ksmall}. Since \(\delta\le7/32\), the slowest power in the
first term is \(a^{-17/32}/L^2\), proving the claim.
\end{proof}

This is absolute normalized accuracy. It gives neither a nonzero
curvature nor a favorable sign for \(\cD_a(S_a)\). To exclude an event,
one would need a lower bound exceeding \(\eta_a\) for the actual
correlated expression.

Two premises of the outer-line shortcut fail here. The formal
shifted-zeta real part approaches \(1/2+c/4<1\), outside the absolutely
convergent series and Euler product used in Section~\ref{sec:fixed};
and
\[
 |C_+/C_-|(T,a+\ii\beta/L)\longrightarrow\ee^{-\beta/2}.
\]
The reflected arm remains present. These facts delimit that shortcut;
they do not exclude other arithmetic comparisons.

\section{Signed phase geometry: disks and exact circles}\label{sec:phases}

\subsection{The full quadratic form}

Suppress the subscript on \(K_a\), and write arbitrary real quadratures
\[
 x_n=d_n\cos\phi_n,\qquad y_n=d_n\sin\phi_n,\qquad
 q(w)=\sum_{n\le N}[x_n\cos(\lambda_nw)-y_n\sin(\lambda_nw)]+K(w).
\]
The actual magnitudes obey \(x_n^2+y_n^2=d_n^2\); a disk relaxation
replaces equality by \(\le\). The two event equations read
\begin{equation}\label{eq:eventquadratures}
 q(0)=\sum_nx_n+K(0)=0,\qquad
 q'(0)=-\sum_n\lambda_n y_n+K'(0)=0.
\end{equation}
The corresponding approximate class allows errors \(O(a^{-3/4})\).
Every exact solution belongs to that approximate class.

At \(\ii\beta\), write \(S=q^\circ\) as \(S(\ii\beta)=X-\ii Y\), where
\begin{align}
 C_n&=\cosh(\beta\lambda_n)-1,&
 B_n&=\sinh(\beta\lambda_n)-\beta\lambda_n,\\
 R_n&=\lambda_n\sinh(\beta\lambda_n),&
 V_n&=\lambda_n(\cosh(\beta\lambda_n)-1),\\
 X&=C\cdot x+\Rea K^\circ(\ii\beta),&
 Y&=B\cdot y-\Ima K^\circ(\ii\beta).
\end{align}
With primes on \(X,Y\) denoting \(\beta\)-derivatives,
\begin{equation}\label{eq:XY}
 \cD_a(S)=\mathfrak d_a(X^2+Y^2)-\beta(XX'+YY').
\end{equation}
Thus, apart from the explicitly retained correction terms, its matrices
are
\[
 Q_x=\mathfrak d_a CC^t-\frac\beta2(CR^t+RC^t),\qquad
 Q_y=\mathfrak d_a BB^t-\frac\beta2(BV^t+VB^t).
\]
These formulas keep all cross terms. They will allow counterexamples
for the full form, independently of any numerical optimization method.
Write \(\cD_2\) for the same functional with coefficient \(2\) in
place of \(\mathfrak d_a\).

\subsection{A two-mode disk construction}

Let
\[
 r=\frac14,\qquad h=\frac{\log2}{2L},\qquad t=\frac{d_2}{2}.
\]
Choose
\[
 x_1=t-K(0),\quad x_2=-t,\quad
 y_1=4K'(0),\quad y_2=0,
\]
and set every other quadrature to zero. Since
\(2^{-7/4}\le t\le2^{-3/2}\) and \(K,K'=O(\chi)\), these points
lie strictly inside their coefficient disks for large \(a\).
Both equations \eqref{eq:eventquadratures} are exact. The resulting
field is
\begin{equation}\label{eq:diskfield}
 q_*(w)=t\{\cos(rw)-\cos((r-h)w)\}+J(w),
 \quad J=K-K(0)\cos(rw)-4K'(0)\sin(rw).
\end{equation}
Its correction \(J\) and fixed derivatives are \(O(\chi)\).

Define the positive number
\begin{equation}\label{eq:Cbeta}
 C_\beta=\frac{(\beta\log2)^2}{4}\sinh^2(\beta/4)
             \left[\frac\beta4\coth(\beta/4)-1\right]>0.
\end{equation}
Positivity follows from \(x\cosh x-\sinh x>0\) for \(x>0\).
For \(U=\cosh(\beta r)-\cosh(\beta(r-h))\), expansion in \(h\) gives
\[
 U=\beta h\sinh(\beta r)+O_\beta(h^2),\qquad
 U(\beta U'-2U)=C_\beta/L^2+O_\beta(L^{-3}).
\]
The uncorrected field is real at \(\ii\beta\), so
\(\cD_2=-t^2U(\beta U'-2U)\). Expanding the entire correction and
carrier contribution yields
\begin{equation}\label{eq:disknegative}
 \cD_a(q_*)=-\frac{t^2C_\beta}{L^2}
   +O_\beta(L^{-3}+\chi/L+\chi^2+a^{-2}L^{-4})<0
\end{equation}
eventually. For example it is at most
\(-2^{-9/2}C_\beta/L^2\) after a finite onset.
This disproves nonnegativity on coefficient disks subject to the two
event equations. Exact magnitudes and the common arithmetic phase law
are absent from the witness.

\subsection{An elementary three-circle lemma}

The next construction restores every magnitude. The following lemma
will also be the local building block for multiplicative phases.

\begin{lemma}[A signed block on three circles]\label{lem:triple}
Let \(b_1>b_2>b_3>0\), \(b_1<b_2+b_3\), and
\(r_1>r_2>r_3\ge0\), with \(r_2>0\). Put \(\rho=r_3/r_2\).
There is a unique \(v\in(0,b_3)\) such that
\[
 A_2=\sqrt{b_2^2-\rho^2v^2},\qquad
 A_3=\sqrt{b_3^2-v^2},\qquad A_2+A_3=b_1.
\]
The quadratures
\begin{equation}\label{eq:triple}
 (x_1,y_1)=(b_1,0),\quad
 (x_2,y_2)=(-A_2,\rho v),\quad
 (x_3,y_3)=(-A_3,-v)
\end{equation}
have exact magnitudes, \(\sum x_i=\sum r_i y_i=0\).
Their field \(p(w)=\sum_i(x_i\cos r_iw-y_i\sin r_iw)\) satisfies
\(p(0)=p'(0)=0\) and \(\cD_2(p)<0\) for every \(\beta>0\).
Any finite sum of these blocks has the same strict sign.
\end{lemma}
\begin{proof}
The sum of the two square roots is continuous and strictly decreasing
on \([0,b_3]\). Its initial value is \(b_2+b_3>b_1\); its final value
is at most \(b_2<b_1\). This proves the unique interior solution.
The circle identities and event equations follow directly.

Writing \(p(\ii\beta)=X-\ii Y\) gives
\begin{align}
 X={}&A_2[\cosh(\beta r_1)-\cosh(\beta r_2)]
      +A_3[\cosh(\beta r_1)-\cosh(\beta r_3)]>0,\\
 Y={}&vr_3\left[\frac{\sinh(\beta r_2)}{r_2}
                       -\frac{\sinh(\beta r_3)}{r_3}\right]\ge0.
 \label{eq:tripleXY}
\end{align}
Use the continuous interpretation at \(r_3=0\). The power series for
\(X\) has positive coefficients in even powers starting with
\(\beta^2\). Consequently \(\beta X'-2X>0\), since the quartic
coefficient is positive. The series for \(Y\) has nonnegative
coefficients in odd powers starting with \(\beta^3\), so
\(\beta Y'-2Y\ge0\).
Equation~\eqref{eq:XY} becomes
\[
 \cD_2(p)=-X(\beta X'-2X)-Y(\beta Y'-2Y)<0.
\]
For sums of blocks the four quantities retain the same signs.
Thus the cross terms between blocks also have the required sign.
\end{proof}

\subsection{All coefficients on their circles}

Partition \(1,\ldots,N\) into consecutive triples, leaving at most two
terminal terms. For a triple \(j,j+1,j+2\), set
\(b_i=d_{j+i-1}\), \(r_i=\lambda_{j+i-1}\).
The logarithmic slope from Lemma~\ref{lem:moments} gives
\[
 b_2+b_3\ge(2^{-3/4}+3^{-3/4})b_1>b_1.
\]
Lemma~\ref{lem:triple} applies, including a possible last frequency
zero. Let \(p\) be the sum of all complete blocks.
Frequency differences in a block are \(O(1/(jL))\).
The event cancellations, applied to the cosine and sine differences
in \eqref{eq:tripleXY} and their analytic versions, give
\[
 \norm{p_j^{(k)}}_{|w|\le R}\le C_R\frac{d_j}{jL}
                 \le \frac{C_R}{Lj^{3/2}},\quad 0\le k\le2.
\]
The possible small \(r_3\) is harmless: use the analytic function
\(\sin(rw)/r\) in the odd term. Summing the convergent \(j^{-3/2}\)
bound yields
\begin{equation}\label{eq:circlefieldbound}
 \norm{p^{(k)}}_{|w|\le R}=O_R(L^{-1}),\qquad 0\le k\le2.
\end{equation}

The head block gives a quantitative negative margin. With
\(\alpha=1-1/\sqrt3>0\), its \(A_2=1-A_3\ge1-d_3\ge\alpha\).
The two first frequencies are \(r=1/4\) and \(r-h\).
All signs in Lemma~\ref{lem:triple} align, so \eqref{eq:Cbeta} implies
\begin{equation}\label{eq:circleleading}
 \cD_2(p)\le-\alpha^2 C_\beta/L^2+O_\beta(L^{-3}).
\end{equation}

Give each leftover terminal term phase zero and denote their sum by
\(b_{\rm tail}(w)\). Since, for \(\ell=0,1\),
\[
 d_{N-\ell}
 =\exp[-(1/4+c/16)L+O(\ee^{-L/2})]=O(\chi),
\]
this tail and its fixed derivatives are \(O_R(\chi)\), and
\(b_{\rm tail}'(0)=0\).
It remains to restore the two event equations in the presence of the
complete \(K\), without leaving a circle.

Vary just the second and third phases of the head block. Their event
map is
\[
 F(\theta_2,\theta_3)=
 \bigl(1+d_2\cos\theta_2+d_3\cos\theta_3,\
       \lambda_2d_2\sin\theta_2+\lambda_3d_3\sin\theta_3\bigr).
\]
At the block solution its Jacobian determinant is
\[
 v\left[\frac{\lambda_3^2}{\lambda_2}A_3+\lambda_2A_2\right]>0.
\]
As \(L\to\infty\), the two frequencies tend to \(1/4\), and
\((1,d_2,d_3)\) tends uniformly to \((1,2^{-s},3^{-s})\) for
\(s\in[5/8,3/4]\). These are a compact family of nondegenerate
triangles, so the inverse Jacobian and the second derivatives of the
map are uniformly bounded. The inverse-function theorem therefore
solves
\[
 F=\bigl(-K(0)-b_{\rm tail}(0),\,K'(0)\bigr)
\]
with phase changes \(O(\chi)\). The complete repaired field
\(q_*=p+E\) satisfies the event equations exactly and
\(\norm{E^{(k)}}_R=O_R(\chi)\) for \(k\le2\).

\begin{theorem}[Exact-circle obstruction]\label{thm:circles}
For each fixed \(\beta>0\), every sufficiently large center in
\eqref{eq:band} admits quadratures with
\(x_n^2+y_n^2=d_n^2\) for every \(n\le N\), the actual correction \(K\),
and \(q_*(0)=q_*'(0)=0\), such that
\begin{equation}\label{eq:circlenegative}
 \cD_a(q_*)\le-\kappa_\beta^{\rm circ}/L^2,\qquad
 \kappa_\beta^{\rm circ}=\tfrac12(1-1/\sqrt3)^2C_\beta>0.
\end{equation}
\end{theorem}
\begin{proof}
Expand \(\cD_a(p+E)\) using \eqref{eq:circlefieldbound}.
The error beyond \eqref{eq:circleleading} is
\[
 O_\beta(\chi/L+\chi^2+a^{-2}L^{-4}).
\]
It is \(o(L^{-2})\), proving the result.
\end{proof}

This exact-circle point is still not multiplicative. The first member
of each block has phase zero, so \(u_1=u_4=1\). The adjusted second
phase remains uniformly in the interior of \((\pi/2,\pi)\). Therefore
\[
 |u_4u_1-u_2^2|=2|\sin\theta_2|\ge\delta_0>0.
\]
Nor is the actual leading carrier phase imposed in this construction.
The next section restores both multiplicativity and that leading phase.

\section{Multiplicativity and a growing calibration budget}\label{sec:multiplicative}

\subsection{The class and the quantitative theorem}

Restore the actual leading phase
\(u_1=\ee^{\ii\Theta(T,a)}\). Let
\(\omega:\N\to\C\) be completely multiplicative, with
\(|\omega(n)|=1\) and \(\omega(1)=1\). Define
\begin{equation}\label{eq:qomega}
 q_\omega(w)=\frac12\sum_{n\le N}d_n
 \left[u_1\omega(n)\ee^{\ii\lambda_nw}
        +\overline{u_1\omega(n)}\ee^{-\ii\lambda_nw}\right]+K(w).
\end{equation}
The actual phase choice is \(\omega_{\can}(n)=n^{-\ii\nu}\).
For arbitrary \(\omega\), all identities
\[
 u_{mn}u_1=u_mu_n,\qquad u_n=u_1\omega(n),
\]
hold, including every retained product \(mn\le N\). The heat weights
are kept exactly; they are not multiplicative:
\begin{equation}\label{eq:weightnonmult}
 \frac{d_{mn}}{d_md_n}
       =\exp\left(\frac{c\log m\log n}{2L}\right).
\end{equation}

Let \(\cQ\) be any set of prescribed primes below \(N/2\), with arbitrary
assigned unit phases. Its cost is
\begin{equation}\label{eq:CQ}
 C_{\cQ}=\prod_{p\in\cQ}(1-p^{-1/2})^{-1}.
\end{equation}
Set
\[
 b_L=\ee^{-(1/4+c/16)L},\qquad
 M=\left\lfloor\frac{N}{40\log N}\right\rfloor,\qquad
 W=M b_L.
\]
The letter \(W\) in this section is a positive population scale; it is
distinct from the three-height quantity \(W_T(a)\).
Eventually
\begin{equation}\label{eq:Wpopulation}
 \frac{\ee^{\delta L}}{80L}\le W\le\frac{\ee^{\delta L}}{10L},
 \qquad \delta=\frac14-\frac c{16}.
\end{equation}
The useful dimensionless error budget is
\begin{equation}\label{eq:calbudget}
 \mathcal E_{\cQ}=\frac{C_{\cQ}L^4}{W}.
\end{equation}

Define fixed positive constants
\begin{align}
 a_0&=\sqrt{17/12}-1,\qquad
 \ell_1=\tfrac12\log(5/3),\qquad \ell_2=\tfrac12\log(3/2),\\
 h_2&=a_0(\ell_1^2-\ell_2^2),\qquad
 h_4=a_0(\ell_1^4-\ell_2^4),\\
 \Gamma_\beta&=\frac{\beta^6h_2h_4}{24},\qquad
 \kappa_\beta=\frac{\Gamma_\beta}{12800}.
 \label{eq:constants}
\end{align}

\begin{theorem}[Uniform multiplicative calibration obstruction]
\label{thm:mult}
Fix \(\beta>0\). Suppose the prescribed prime sets satisfy
\(\mathcal E_{\cQ}\to0\) uniformly in \eqref{eq:band}. At every
sufficiently large center there is a completely multiplicative
\(\omega\) attaining all prescribed phases such that the full
\eqref{eq:qomega}, with the actual \(u_1,d_n,K\), satisfies
\begin{equation}\label{eq:multnegative}
 q_\omega(0)=q_\omega'(0)=0,\qquad
 \cD_a(q_\omega)\le
 -\kappa_\beta\frac{\ee^{(1/2-c/8)L}}{L^8}.
\end{equation}
The conclusion holds for any prescribed phases, not only canonical
ones. Its onset is finite and uninstantiated.
\end{theorem}

The proof combines three ingredients: isolated primes that can be
adjusted independently without breaking multiplicativity; a bounded
complete complement obtained by finite character orthogonality; and
a uniform two-parameter correction of the event equations.
The general isolated-prime first-jet mechanism predates RSCT
\cite{collision}. The extension here is its negative signed margin,
complete correction, explicit calibration budget, and simultaneous
fourth-derivative compatibility in Section~\ref{sec:mirror}.

\subsection{A reserved prime population}

Use the intervals
\[
 I_1=[11/20,3/5],\quad I_2=[7/10,3/4],\quad
 I_3=[17/20,9/10].
\]
The prime number theorem \cite{dlmfpnt} gives
\[
 \#\{p:p/N\in I_j\}=(1/20+o(1))N/\log N.
\]
Select \(M\) primes from each interval and arrange them in \(M\)
triples \(p_{j,1}<p_{j,2}<p_{j,3}\). Write \(\cR\) for their union.
Every reserved prime exceeds \(N/2\); among \(n\le N\), its only
multiple is itself. Changing its phase therefore changes no
complementary term. In particular \(\cR\cap\cQ=\varnothing\).

For a reserved prime write \(v=L/2-\log p\). Direct substitution gives
\begin{equation}\label{eq:primeweights}
 d_p=b_L\ee^{v/2+cv^2/(4L)},\qquad \lambda_p=v/(2L).
\end{equation}
Thus \(d_p/b_L\) and \(L\lambda_p\) lie in fixed positive compact
intervals, separated between the three groups. The floor in \(M\)
and \(N\sim\ee^{L/2}\) give \eqref{eq:Wpopulation}. For instance,
\(M\ge N/(80\log N)\), \(N\ge\ee^{L/2}/2\), and
\(\log N\le L/2\) give its lower bound.
Only primes in fixed relative intervals have been used; no
short-interval theorem or effective numerical onset is invoked.

\subsection{The complete complement and its dependence on calibration}

Keep the prescribed phases on \(\cQ\). Assign independent uniform circle
phases to every other nonreserved prime at most \(N\). Let \(B(w)\)
be the complete sum in \eqref{eq:qomega} over \(n\notin\cR\), excluding
only the separate correction \(K\). It includes \(n=1\), every
unreserved prime, and every composite through \(N\).

\begin{lemma}[Finite orthogonality with prescribed primes]\label{lem:complement}
For any one of the observations
\[
 B^{(j)}(0)\quad(0\le j\le4),\qquad B(\ii\beta),\quad B'(\ii\beta),
\]
the mean square over the free prime phases is at most
\(9\ee^{\beta/2}C_{\cQ}^2\).
There is consequently one choice of the complementary phases for which
the sum of all seven squared observations is at most
\(63\ee^{\beta/2}C_{\cQ}^2\). If only orders zero and one at zero
and \(\ii\beta\) are requested, the corresponding constant is
\(36\ee^{\beta/2}C_{\cQ}^2\).
\end{lemma}
\begin{proof}
The elementary weight bounds are
\begin{equation}\label{eq:weightbounds}
 d_n^2\le n^{-1-c/4}\le n^{-9/8},\qquad
 \sum_{n\le N}d_n^2\le9,\qquad d_{rm}\le r^{-1/2}d_m
 \quad(rm\le N).
\end{equation}
The last follows by integrating the logarithmic slope bound
\(\dd\log d(x)/\dd\log x\le-1/2\).
Every \(n\) has a unique factorization \(n=rm\), with \(r\)
supported on \(\cQ\) and \(m\) free of primes in \(\cQ\).
Group the terms in either arm by the character \(\omega(m)\) of the
free prime variables. Distinct \(m\)'s give distinct characters.
No reserved prime enters a complementary character because its only
retained multiple was removed.

For each observation the absolute value of a character coefficient is
at most
\[
 \ee^{\beta/4}d_m
       \sum_{r:\,\cQ\text{-supported}}r^{-1/2}
 =\ee^{\beta/4}C_{\cQ}d_m;
\]
all fixed derivative factors obey \(\lambda_n^j\le1\).
Integration of \(\ee^{\ii k\theta}\) annihilates unequal characters,
giving the bound by \eqref{eq:weightbounds}. The deterministic
\(m=1\) group is included. The reflected arm has the same estimate;
\(|(A_++A_-)/2|^2\le(|A_+|^2+|A_-|^2)/2\) pays their cross term.
Sum the expectations and select a point no larger than the mean.
\end{proof}

This is finite integration on an auxiliary phase space. It is neither
a numerical experiment nor an average over actual heights. In
particular, the selected complement has controlled values at the
listed observations; we have not asserted a uniform disk bound for it.

\subsection{Negative blocks of isolated primes}

For each reserved triple apply Lemma~\ref{lem:triple} with
\[
 b_i=d_{p_{j,i}},\qquad r_i=\lambda_{p_{j,i}}.
\]
Its triangle condition holds with a uniform margin:
\[
 b_1/b_3\le(p_{j,3}/p_{j,1})^{3/4}\le(18/11)^{3/4}<2,
 \qquad b_2>b_3.
\]
After division by \(b_L\) and multiplication of frequencies by \(L\),
the block parameters form a compact nondegenerate family. Thus
\(A_2,A_3,v,\rho v\asymp b_L\).
Also
\[
 A_2\ge b_1-b_3\ge a_0b_L,\qquad
 Lr_1\ge\ell_1,\quad Lr_2\le\ell_2
\]
for large \(N\). For the last upper bound,
\(\ee^{L/2}/N\le21/20\) when \(N\ge20\).

Let \(p\) be the sum of these blocks. The positive power-series
coefficients in \eqref{eq:tripleXY} give
\begin{equation}\label{eq:primeXG}
 X\ge\frac{\beta^2h_2}{2}\frac W{L^2},\qquad
 \beta X'-2X\ge\frac{\beta^4h_4}{12}\frac W{L^4}.
\end{equation}
The first bound uses the quadratic cosh term; the second uses the
quartic term, whose factor under \(\beta\partial_\beta-2\) is two.
All blocks have aligned signs. Hence
\begin{equation}\label{eq:primeD}
 \cD_2(p)\le-\Gamma_\beta W^2/L^6,\qquad
 |p(\ii\beta)|+|p'(\ii\beta)|=O_\beta(W/L^2).
\end{equation}
At zero the same block algebra gives the derivatives needed later:
\begin{equation}\label{eq:primejets}
 p''(0)\le-h_2W/L^2,\qquad
 p^{(4)}(0)\ge h_4W/L^4.
\end{equation}
Indeed \(p''(0)=-\sum r_i^2x_i\) and
\(p^{(4)}(0)=\sum r_i^4x_i\); the \(A_2\) terms alone provide these
two bounds, and the \(A_3\) terms have the same signs.

\subsection{Uniform repair of the exact event}

In each reserved triple shift the phase of its second member by the
same \(t\), and its third by the same \(s\). The normalized event map is
\[
 F(t,s)=\left(W^{-1}\sum_{\cR}x_n(t,s),\
                   LW^{-1}\sum_{\cR}\lambda_n y_n(t,s)\right).
\]
At the block configuration \(F(0,0)=0\), and
\begin{equation}\label{eq:repairJacobian}
 DF(0)=
 \frac1W\begin{pmatrix}
 -\sum y_2&-\sum y_3\\
 L\sum r_2x_2&L\sum r_3x_3
 \end{pmatrix}
 =\begin{pmatrix}-A_J&B_J\\-C_J&-D_J\end{pmatrix},
 \qquad A_J,B_J,C_J,D_J\asymp1.
\end{equation}
The signs come from \(x_2,x_3<0\) and \(y_2>0>y_3\).
Thus \(\det DF=A_JD_J+B_JC_J\) is uniformly positive.
The inverse and the second derivatives on a fixed small parameter disk
are uniformly bounded.

Use Lemma~\ref{lem:complement} to fix a complementary phase choice.
Solve
\begin{equation}\label{eq:repairtarget}
 F(t,s)=\left(-\frac{B(0)+K(0)}W,\
                    \frac{L(B'(0)+K'(0))}W\right).
\end{equation}
The target has norm \(O_\beta(C_{\cQ}L/W)\), which tends to zero under
\eqref{eq:calbudget}. A quantitative inverse-function argument gives
\begin{equation}\label{eq:repairsize}
 |t|+|s|=O_\beta(C_{\cQ}L/W).
\end{equation}
Explicitly, apply the inverse of \eqref{eq:repairJacobian} to the
target and to the quadratic Taylor remainder. Uniformly bounded
second derivatives make this a contraction on a ball of radius a
fixed multiple of the target norm, once that norm is small.

Assign the final prime value \(\omega(p)=\ee^{\ii\theta_p}/u_1\)
at each reserved prime, keep every complementary prime value, and
extend completely multiplicatively. Isolation ensures that no
complementary coefficient changes. If \(\Delta p\) is the reserved
repair, the complete field
\[
 q_*=B+p+\Delta p+K
\]
satisfies \(q_*(0)=q_*'(0)=0\) exactly.

The repair must be estimated after the event equations have been used.
Since \(p=p^\circ\) and \(q_*=q_*^\circ\),
\[
 q_*=p+E,\qquad E=B^\circ+K^\circ+(\Delta p)^\circ.
\]
For a changed mode the subtracted value and derivative on a fixed disk
are \(O_\beta(d_n(|t|+|s|)\lambda_n^2)\).
Equations~\eqref{eq:primeweights} and~\eqref{eq:repairsize} therefore give
\[
 |(\Delta p)^\circ(\ii\beta)|
       +|((\Delta p)^\circ)'(\ii\beta)|
       =O_\beta(C_{\cQ}/L).
\]
The complete background observations from Lemma~\ref{lem:complement}
and the correction bound give
\begin{equation}\label{eq:Ebound}
 |E(\ii\beta)|+|E'(\ii\beta)|=O_\beta(C_{\cQ}).
\end{equation}
Raw repair values could be \(O(C_{\cQ}L)\); event subtraction is what
avoids that unnecessary loss.

Expanding the full quadratic expression, including background and
correction cross terms and \(\mathfrak d_a-2\), now yields
\begin{equation}\label{eq:fullprimeerror}
 \cD_a(q_*)\le-\Gamma_\beta\frac{W^2}{L^6}
 +O_\beta\left(\frac{C_{\cQ}W}{L^2}+C_{\cQ}^2
                         +a^{-2}\frac{W^2}{L^6}\right).
\end{equation}
The relative errors are respectively
\(O_\beta(\mathcal E_{\cQ})\),
\(O_\beta(\mathcal E_{\cQ}^2/L^2)\), and \(O_\beta(a^{-2})\).
After they total at most \(\Gamma_\beta/2\), the lower bound in
\eqref{eq:Wpopulation} proves
\[
 \cD_a(q_*)\le-\tfrac12\Gamma_\beta W^2/L^6
 \le-\frac{\Gamma_\beta}{12800}
                         \frac{\ee^{(1/2-c/8)L}}{L^8}.
\]
This completes the proof of Theorem~\ref{thm:mult}.
The small physical approximation error from
Proposition~\ref{prop:Derror} is not assigned to this artificial field;
no twisted physical approximation theorem has been asserted.

\subsection{Fixed and growing actual-prime calibration}

For a fixed finite set \(\mathcal P_0\) and a fixed prime
\(q_0\notin\mathcal P_0\), prescribe
\[
 \omega(p)=p^{-\ii\nu}\ (p\in\mathcal P_0),\qquad
 \omega(q_0)=-q_0^{-\ii\nu}.
\]
Here \(C_{\cQ}\) is fixed and \(\mathcal E_{\cQ}\to0\).
Theorem~\ref{thm:mult} recovers the fixed-calibration result, with the
explicit same-height defect
\begin{equation}\label{eq:phasedefect}
 |u_{q_0}-u_{q_0}^{\can}|=2.
\end{equation}

The displayed budget also permits a growing set.
\begin{corollary}[A sufficient growing calibration range]\label{cor:growing}
Fix \(0<\gamma<9/256\), and put
\begin{equation}\label{eq:Y}
 Y=\gamma L^2(\log L)^2.
\end{equation}
For every sufficiently large center, all actual prime phases
\(\omega(p)=p^{-\ii\nu}\) for \(p\le Y\) can be retained in
Theorem~\ref{thm:mult}. In addition one can choose a prime
\(q_0\in(Y,2Y]\) with the opposite actual phase, so
\eqref{eq:phasedefect} holds.
\end{corollary}
\begin{proof}
Partial summation and the prime number theorem give
\[
 \sum_{p\le Y}p^{-1/2}
 =\frac{\pi(Y)}{\sqrt Y}
       +\frac12\int_2^Y\frac{\pi(t)}{t^{3/2}}\dd t
 \sim\frac{2\sqrt Y}{\log Y}.
\]
Expanding \(-\log(1-p^{-1/2})\), the terms of order at least two
total \(O(\sum_{p\le Y}p^{-1})=O(\log Y)\), already small enough here.
Thus
\[
 \log\prod_{p\le Y}(1-p^{-1/2})^{-1}
 =\frac{2\sqrt Y}{\log Y}(1+o(1))
 =(\sqrt\gamma+o(1))L.
\]
Since \(\delta\ge3/16\), \eqref{eq:Wpopulation} implies
\[
 \mathcal E_{\cQ}\le80L^5
       \exp[-(3/16-\sqrt\gamma+o(1))L]\longrightarrow0.
\]
The prime number theorem supplies a \(q_0\in(Y,2Y]\); it is below
\(N/2\) for large \(L\). Its additional cost
\((1-q_0^{-1/2})^{-1}\) does not change the budget.
\end{proof}

For example \(\gamma=1/64\) calibrates
\((1+o(1))L^2\log L/128\) primes. Complete multiplicativity then fixes
every retained \(Y\)-smooth integer at its actual phase.
The coefficient \(9/256\) is sufficient for this proof, not a claimed
sharp boundary. The explicit defect above proves failure of the actual
same-height phase law. Realizing a similar phase pattern at another
height would change \(d_n,N,u_1,K\) and the time-height relation; no
such realization has been proved.

\section{A forced-mirror fourth-derivative condition}\label{sec:mirror}

\subsection{Deflation and the exact physical inequality}

At an all-real event \(a>0\) of multiplicity \(m\ge2\), deflate two
copies at \(a\):
\[
 G_a(z)=\frac{H_T(z)}{(z-a)^2}.
\]
The remaining zero multiset still contains at least two copies at
\(-a\). For real \(z\) away from the zeros, logarithmic differentiation
of its product gives
\[
 (G_a'(z))^2-G_a(z)G_a''(z)
  =G_a(z)^2\sum_{\rho:\,G_a(\rho)=0}(z-\rho)^{-2}
  \ge\frac{2G_a(z)^2}{(z+a)^2}.
\]
The inequality extends to \(z=a\) by continuity, including when \(G_a(a)=0\).
Write \(H_j=H_T^{(j)}(a)\). Taylor expansion gives
\[
 G_a(a)=H_2/2,\qquad G_a'(a)=H_3/6,\qquad G_a''(a)=H_4/12.
\]
Multiplication by \(72\) proves the inherited mirror condition
\begin{equation}\label{eq:physicaljet}
 \boxed{2H_3^2-3H_2H_4\ge\frac9{a^2}H_2^2.}
\end{equation}
For \(m\ge4\), both sides vanish; for \(m=3\) the inequality remains
valid with \(H_2=0\). No restriction to exact doubles was made.
The deflation condition and its earlier model compatibility are owned
by \cite{collision}; this section develops its RSCT normalization and
joint calibrated use.

\subsection{Cancellation of the carrier derivatives}

Let \(f_j=\partial_w^jf_a(0)\), with \(f_0=f_1=0\), and evaluate
\(b=A_z/A\) and \(b_z\) at \(z=a\), holding \(T\) fixed.
Differentiating \(H_T=A f_a\) gives
\begin{align}
 H_2/A&=L^2f_2,\\
 H_3/A&=L^3f_3+3bL^2f_2,\\
 H_4/A&=L^4f_4+4bL^3f_3+6(b^2+b_z)L^2f_2.
\end{align}
Substitute in \eqref{eq:physicaljet}. The \(bf_2f_3\) and \(b^2f_2^2\)
terms cancel exactly. After division by \(A(T,a)^2L^6>0\), the condition is
\begin{equation}\label{eq:Jdef}
 \cJ_a(f_a):=2f_3^2-3f_2f_4-\mu_a f_2^2\ge0,
 \qquad \mu_a=\frac{18b_z(a)+9/a^2}{L^2}.
\end{equation}
From \eqref{eq:bcarrier},
\[
 b_z(a)=-\frac7{4a^2}+\frac{T(1-L)}{8a^2},
\]
and hence
\begin{equation}\label{eq:mu}
 \mu_a=\frac{-45/2-9c/4+9c/(4L)}{a^2L^2}<0\qquad(L\ge1).
\end{equation}
The normalization therefore retains a small positive contribution
\(-\mu_a f_2^2\) in the mirror functional. A carrier derivative has
not been discarded as negligible before the algebraic cancellation.

\subsection{Payment for the fourth derivative}

Theorem~\ref{thm:local} uses an analytic error on a larger disk, so
there are explicit Cauchy payments
\[
 |f_j-q_j|\le\epsilon_j,\qquad
 \epsilon_j=j!R^{-j}C_{2R}a^{-3/4},\qquad 2\le j\le4,
\]
where \(q_j=\partial_w^jq_a(0)\).
All these jets are real. Expanding \eqref{eq:Jdef} gives
\begin{equation}\label{eq:jeterror}
\begin{split}
 |\cJ_a(f_a)-\cJ_a(q_a)|\le\Delta_{\rm jet}:={}&
 4|q_3|\epsilon_3+2\epsilon_3^2
 +3|q_2|\epsilon_4+3|q_4|\epsilon_2+3\epsilon_2\epsilon_4\\
 &+|\mu_a|(2|q_2|\epsilon_2+\epsilon_2^2).
\end{split}
\end{equation}
The moment estimate and the correction imply
\(|q_j|\le M_j+O(\chi)\). Therefore
\begin{equation}\label{eq:jeterrorsize}
 \Delta_{\rm jet}
 =O\left(a^{-3/4}\ee^{\delta L}/L^2+a^{-3/2}\right)=o(1).
\end{equation}
For the actual approximate field, the justified necessary statement is
\begin{equation}\label{eq:approxjetcondition}
 |q_a(0)|,\ |q_a'(0)|=O(a^{-3/4}),\qquad
 \cJ_a(q_a)\ge-\Delta_{\rm jet}.
\end{equation}
Exact nonnegativity of the approximate functional is not silently
substituted for this paid inequality.

\subsection{One point satisfies both signed tests}

\begin{theorem}[Joint mirror compatibility]\label{thm:joint}
Under the hypotheses of Theorem~\ref{thm:mult}, the same
multiplicative point can be chosen to satisfy
\begin{align}
 q_*(0)&=q_*'(0)=0,\qquad q_*''(0)<0,\qquad q_*^{(4)}(0)>0,\\
 \cD_a(q_*)&\le-\kappa_\beta\frac{\ee^{(1/2-c/8)L}}{L^8},
 \label{eq:jointD}\\
 \cJ_a(q_*)&\ge c_J\frac{\ee^{(1/2-c/8)L}}{L^8}>0,\qquad
 c_J=\frac{3h_2h_4}{25600}.
 \label{eq:jointJ}
\end{align}
In particular, Corollary~\ref{cor:growing}, including the opposite phase
at a prime \(q_0\in(Y,2Y]\), remains valid simultaneously.
\end{theorem}
\begin{proof}
Make the seven-observation choice in Lemma~\ref{lem:complement}, rather
than only the four-observation choice. The event repair and signed
estimate in Section~\ref{sec:multiplicative} are unchanged, with
constants still independent of \(\cQ\).

For each \(j=2,3,4\), the reserved phase repair has derivative cost
\[
 |(\Delta p)^{(j)}(0)|
 \le C W(|t|+|s|)L^{-j}
 =O_\beta(C_{\cQ}L^{1-j}).
\]
The chosen complementary derivatives are \(O_\beta(C_{\cQ})\);
the entire correction has derivatives \(O(\chi)\).
The budget \(\mathcal E_{\cQ}\to0\) implies
\(C_{\cQ}=o(W/L^4)\). Thus the errors are smaller than both leading
quantities in \eqref{eq:primejets}, and eventually
\begin{equation}\label{eq:finaljets}
 q_*''(0)\le-\tfrac12h_2W/L^2,\qquad
 q_*^{(4)}(0)\ge\tfrac12h_4W/L^4.
\end{equation}
All jets are real, and \(\mu_a<0\). In \eqref{eq:Jdef},
the \(2q_3^2\) and \(-\mu_a q_2^2\) terms are nonnegative.
The cross term therefore gives
\[
 \cJ_a(q_*)\ge-3q_*''(0)q_*^{(4)}(0)
        \ge\tfrac34h_2h_4W^2/L^6.
\]
Use \eqref{eq:Wpopulation} to obtain \eqref{eq:jointJ}.
Equation~\eqref{eq:jointD} holds for this very same point by the
previous signed-error calculation.
\end{proof}

The model now has an exact local double zero at \(w=0\), since its
second derivative is nonzero. It passes the exact mirror functional,
which is stronger than merely passing the approximate necessary
condition \eqref{eq:approxjetcondition}. Nevertheless this is not
a physical collision: the model is not identified with a globally
even canonical field, and no all-real global zero configuration is
constructed. The opposite calibrated-prime defect explicitly forbids
the actual same-height phase law.

The conclusion is limited to the one mirror inequality tested here.
It does not establish compatibility with every finite collection of
Laguerre inequalities, every higher event jet, or global all-realness.
Each stronger constraint would need its own derivation, approximation
cost, and simultaneous construction or exclusion argument.

\section{Controls, ownership, and evidence}\label{sec:controls}

\subsection{Exact heat-flow controls}

The signs above use all-realness, evenness, and the complete zero
population. Those properties alone cannot yield the opposing
arithmetic sign. Two elementary controls make this precise.

For \(a,T>0\), put \(s=t-T\) and
\begin{equation}\label{eq:polynomial}
 P(t,z)=(z^2-a^2)^2-s(12z^2-4a^2)+12s^2.
\end{equation}
It satisfies \(P_t=-P_{zz}\). Its squared roots are
\[
 a^2+6s\pm\sqrt{8a^2s+24s^2}.
\]
For \(s>0\) they are distinct positive numbers: their sum is positive
and their product is \(a^4+4a^2s+12s^2>0\).
For \(-a^2/3<s<0\) they are nonreal; for \(s\le-a^2/3\) they are
negative. At \(s=0\), \(\pm a\) are double. Thus \(T\) is its exact
all-real threshold. The arm \(\mathscr A=P/2+\ii d\), \(d\ne0\) real,
has a stationary phase at the event; \(d=0\) gives an arm zero.
At the event its Poisson expression has \(Q=2k\) and \(\cB=0\),
because the entire zero population is the reflected double pair.
It fails the canonical theta arithmetic and infinite zero population.

An infinite positive-source control is more relevant to strict
negativity. Fix \(\ell=2\pi/3\), and for each source parameter \(T>0\)
define
\begin{equation}\label{eq:atomic}
 G_T(t,z)=\frac49\left[\frac34+\ee^{\ell^2(t-T)}\cos(\ell z)
                 +\frac12\ee^{4\ell^2(t-T)}\cos(2\ell z)\right].
\end{equation}
Its half-line cosine source has three positive masses, at \(0,\ell,2\ell\),
and it satisfies the same backward heat equation. At \(t=T\),
\[
 G_T(T,z)=\frac49(\cos(\ell z)+1/2)^2.
\]
For completeness, set \(r=\ee^{\ell^2(t-T)}\) and \(x=\cos(\ell z)\).
The bracket is
\[
 r^4x^2+rx+\frac34-\frac12r^4.
\]
Its discriminant is \(r^2(1-3r^2+2r^6)\).
For \(r>1\) the discriminant is positive, the vertex is in \((-1,0)\),
the vertex value is negative, and both endpoint values at \(x=\pm1\)
are positive. Both distinct roots are in \((-1,1)\), so all zeros
are simple and real. For \(r<1\), either the discriminant is negative,
or \(r^2\le(\sqrt3-1)/2<1/2\). In the latter case the root sum
\(-r^{-3}<-2\) forces a root below \(-1\). Thus nonreal zeros occur
for every \(t<T\). This is an exact positive-source threshold.

The centers \(a_j=1+3j\) are double zeros at \(T\).
For large \(a_j\), choose the source parameter
\(T=c/\log(a_j/(4\pi))\). This tests the same time-height band using
different members of the source family; it does not give a fixed
source with changing threshold. Trigonometric factorization gives
\[
 \log|G_T(T,a_j+\ii y)|
 =\log(1/9)+2\ell y+
             \log(1-\ee^{-\ell y})+\log(1-\ee^{-3\ell y}).
\]
The affine part cancels from the three-height statistic.
At the four remaining negative-log terms at \(y=4,8\),
\(-\log(1-\ee^{-by})<1/255\) for \(b=\ell,3\ell\).
The two middle log terms are negative. Hence
\begin{equation}\label{eq:controlcap}
 W_{G_T}(a_j)<4/255<1/6<13/72.
\end{equation}
The same event decomposition as \eqref{eq:Wbackground} therefore
permits the negative background required by the scalar cap.
Its infinitely many additional real zeros also give \(\cB<0\) at
every event observation \(y>0\), including \(y=\beta/L\).
This control fails unchanged theta coefficients, full gamma/profile
growth, and eventual simplicity at its own threshold. The source is
the BW-01 control cited in the original RSCT notes, not a canonical
theta construction.

\subsection{What each negative result excludes}

\begin{center}
\small
\begin{tabularx}{\textwidth}{@{}>{\raggedright\arraybackslash}p{0.21\textwidth}
 >{\raggedright\arraybackslash}X >{\raggedright\arraybackslash}X@{}}
\toprule
Tested information & Quantitative conclusion & Missing canonical premise\\
\midrule
Fixed heights \(4,6,8\)&
\(\mathcal R_{\rm bg}<-1/36\) under an event; matrix cap
\((6/13)L+O(1)\)&
A favorable arithmetic estimate for the signed background\\
\addlinespace
Coefficient disks&
Negative value of order \(L^{-2}\), exact first two event equations&
Exact magnitudes and common phase law\\
\addlinespace
All coefficient circles&
Negative value of order \(L^{-2}\), full correction&
Multiplicativity and actual leading phase\\
\addlinespace
Multiplicative unit phases&
Negative value of order \(\ee^{(1/2-c/8)L}/L^8\)&
Actual remaining phases at the same height\\
\addlinespace
Actual primes \(p\le\gamma L^2(\log L)^2\)&
Same negative margin for \(0<\gamma<9/256\)&
Uncalibrated actual phases and global canonical field\\
\addlinespace
The preceding class plus the mirror condition&
One point has \(\cD_a<0\) and \(\cJ_a>0\), with explicit margins&
Global all-realness, further constraints, actual phase trajectory\\
\bottomrule
\end{tabularx}
\end{center}

The phase constructions are counterexamples to a proposed
nonnegative lower bound over the displayed estimating classes.
They do not supply a canonical negative value or a physical
collision. The complete approximation theorem only applies to the
canonical field; it is not transferred to a twisted point.

The broader
\repo{result-summaries/12_collision_controls_and_screening.md}{collision-control catalog}
was screened analytically on its recorded witness domains.
Euler-factor and translated controls alter the canonical coefficients;
their nonreal witnesses are not automatically all-real threshold events.
Weighted-square, finite-tail, and finite-prime-orbit controls fail the
full source or reflection; finite-tail controls have no all-real time.
Reflected super-Gaussian, smooth-shape, and positive-primitive controls
fail theta square-series arithmetic. Atomic, finite-profile, and compact
continuous threshold constructions fail the complete canonical profile.
Conditional endpoint controls retain their conditional time-zero premise
and fail ordinary Dirichlet arithmetic. Variable-readout and heat
counter-kernel controls alter the readout, source, or physical generator.
Relaxed zero/profile populations do not identify the exact canonical
integer zero multiset and full boundary field.
No control is moved to a witness domain not proved for it.

\subsection{Source crosswalk and numerical role}\label{sec:crosswalk}

The proofs above reorganize the source notes by mathematical dependence.
The following crosswalk preserves their original ownership:

\begin{center}
\small
\begin{tabularx}{\textwidth}{@{}>{\raggedright\arraybackslash}p{0.20\textwidth}
 >{\raggedright\arraybackslash}X >{\raggedright\arraybackslash}p{0.15\textwidth}@{}}
\toprule
Source & Result carried into this manuscript & Location\\
\midrule
\cite{collision,polymath}&Finite attainment, local splitting, paired
product, prior absolute-population and multiplicative mechanisms,
mirror deflation&\S\S\ref{sec:geometry}, \ref{sec:multiplicative}, \ref{sec:mirror}\\
\addlinespace
\cite{hda,gnpe01,s04}&Buffered complete field provider, gamma correction,
outer comparison, prior local approximation and terminal scale&
\S\S\ref{sec:representation}, \ref{sec:fixed}, \ref{sec:micro}\\
\addlinespace
RSCT R1, Note03 \cite{r1}&Arm alternatives, reflected Poisson expression,
complete error and small-\(y\) cost&\S\S\ref{sec:geometry}, \ref{sec:representation}\\
\addlinespace
R2, Note04 \cite{r2}&Signed three-height series, noncentral requirement,
kernel obstruction, threshold control&\S\ref{sec:fixed}, \S\ref{sec:controls}\\
\addlinespace
R3, Note05 \cite{r3}&Complete derivative acquisition and matrix cap&
\S\ref{sec:fixed}\\
\addlinespace
R4, Note06 \cite{r4}&Microscopic model, fixed-order moments, event
subtraction and signed-error payment&\S\ref{sec:micro}\\
\addlinespace
R5--R6, Notes07--08 \cite{r5,r6}&Full quadratic form, disk and
exact-circle constructions&\S\ref{sec:phases}\\
\addlinespace
R7--R8, Notes09--10 \cite{r7,r8}&Signed multiplicative extension and
uniform growing calibration budget&\S\ref{sec:multiplicative}\\
\addlinespace
R9, Note11 \cite{r9}&Normalized mirror condition, fourth-derivative error,
joint construction&\S\ref{sec:mirror}\\
\bottomrule
\end{tabularx}
\end{center}

R1 has the source status \emph{elementary identities and a scoped failed
acquisition}. R2--R6 are \emph{proved in-note from named inherited
approximation inputs; same-assistant checking only}. R7--R9 are
\emph{proved in-note from named classical and inherited approximation inputs;
same-assistant checking only}. Their associated same-assistant reviews
remain linked from the source notes and the manuscript source catalog.
In particular, the present exposition is not a new independent review
of HDA CS3/CS6.

No RSCT numerical experiment is required for the results presented here.
The rational inequalities, integral evaluations, finite character
averages, and asymptotic estimates above are analytical arguments.
The small records in the investigation's \texttt{numerics/} folder
principally document file identities, links, and execution hygiene.
Those records and the manuscript build checks do not validate the
mathematical claims. Inherited external numerical bounds are not used
to strengthen this manuscript's analytical conclusions. All onsets
depending on the internal approximation constants remain uninstantiated.

\section{Outstanding tasks}\label{sec:outstanding}

\subsection{Obtain a sign from the actual remaining arithmetic}

The central unpaid step is a useful lower bound for
\(\cD_a(q_{\can}^\circ)\) under the actual event equations,
exceeding the complete error \(\eta_a\) in \eqref{eq:eta}.
The phase law is
\[
 u_n^{\can}=\ee^{\ii\Theta(T,a)}
            \ee^{-\ii(a/2+\pi T/8)\log n}.
\]
Every phase, weight, cutoff, carrier, and correction comes from the same
\((a,T)\). Multiplicativity alone, even supplemented by the growing
calibration of Corollary~\ref{cor:growing} and the one mirror inequality,
does not enforce the desired lower bound.
A further argument must use information absent from
Theorem~\ref{thm:joint}, or select a materially different measurement.
Merely evaluating the field more accurately does not supply its sign.

The next bounded analytical question should specify the new information,
the inequality it is meant to prove, the event domain, and the error
budget. This manuscript does not identify an already working method for
that question. RSCT-E01 remains reserved rather than an executable proof
task; the conceptual investigation remains open.

\subsection{Test stronger event information jointly}

Full all-realness contains more information than
\eqref{eq:physicaljet}. Higher deflated Laguerre inequalities, additional
event jets for \(m>2\), or constraints on a complete noncentral zero
measure could supply further necessary conditions.
Their utility is not settled by the present constructions.
For each proposed addition one must derive the physical condition,
transform the carrier exactly, pay the required derivatives, and test
the enlarged set of conditions at one common point.
Separate compatibility witnesses do not prove joint compatibility.
Conversely, the one successful joint construction here does not imply
that every finite family of conditions is ineffective.

\subsection{Distinguish useful mean estimates from exceptional events}

The neighboring Green-potential program has positive results for an
actual canonical observable: at \(T=1/L_X\), its dyadic moments tend to
zero for every fixed positive power, with corresponding subcritical
exponential-moment conclusions \cite{gnp19}. These are moving-time
average statements, with no present rate or every-center conclusion.
At an all-real threshold its Green value is already zero, even at a
multiple real zero \cite{gnp20}; the value alone cannot exclude the
threshold event.

Three tools from this investigation are suitable for discussion there:
the explicit dependence on a growing calibrated prime set; the complete
local approximation and derivative-error propagation; and the
calibrated phase constructions as tests of proposed estimates.
Their original measurements and quantifiers must be retained during a
transfer. In particular, the auxiliary phase average in
Lemma~\ref{lem:complement} is not an actual-height average.

A concrete proposed question is to quantify the Green program's
fixed-cutoff residual mean
\[
 \frac1U\int_U^{2U}R_P(X,1)\dd X=\mu_P+E(U,P),
 \qquad
 \mu_P=\zeta(3/2)\prod_{p\le P}(1-p^{-3/2})-1,
\]
uniformly for a declared growing range of \(P\), ideally including a
cutoff of order \((\log U)^2\). Here \(R_P\) is the complete normalized
two-arm residual of the Green program, not \(\cD_a\).
This is an unproved quantitative question. A cutoff fixed on each
dyadic block but growing between blocks is an intermediate target;
passing to a cutoff varying with \(X\) inside the integral requires
another argument. Even a successful mean estimate leaves sparse
exceptional centers.
\repo{papers/investigations-5/09_green_potentials_newman_flow/01_riemann_siegel_evolution/brainstorm/external/01_RSCT_MANUSCRIPT_COMPARISON_AND_PROPOSED_QUESTIONS_20261112.md}{The companion comparison note}
gives the exact residual definition,
source links, and the additional obligations for a local deflation
transfer. No such transfer is claimed proved here.

\subsection{Cover the full event domain}

Even a favorable result throughout \eqref{eq:band} would cover only a
specified asymptotic regime. The original threshold problem includes
bounded centers, other values of \(T\log(a/(4\pi))\), and arbitrary
multiplicity. Known fixed-positive-time high-zero reality and simplicity
are inherited inputs, not new coverage from this investigation.
A complete argument must state the remaining domain partition and
prove all required overlaps without applying fixed-geometry estimates
to changing geometries.

\subsection{Audit inputs and make constants effective where useful}

The field provider, its larger complex charts, Gaussian exterior, and
fractional-cutoff correction need an independent specialist audit before
the present conditional proof chain is treated as independently
established. The local and matrix constants are finite, but their
numerical values and onsets have not been certified. Making them
effective is useful only for a clearly specified application with an
available arithmetic margin.
Finite computations could then support that application by checking
normalizations, evaluating complete expressions with errors, or covering
a bounded remaining region. They cannot replace the missing
all-center arithmetic argument or validate it by sampled agreement.

\begin{thebibliography}{99}
\addcontentsline{toc}{section}{References and analytical sources}

\bibitem{polymath}
D. H. J. Polymath,
\emph{Effective approximation of heat flow evolution of the Riemann
\(\xi\) function, and a new upper bound for the de Bruijn--Newman constant},
arXiv:1904.12438v2, 4 August 2019.
\href{https://arxiv.org/html/1904.12438v2}{Primary text}.
Theorem 1.5, Proposition 3.1 and Section 4, equations (35)--(39),
are the passages used here; rechecked 12 November 2026.

\bibitem{dlmfzeta}
NIST Digital Library of Mathematical Functions,
\href{https://dlmf.nist.gov/25.2}{Section 25.2, Definition and Expansions}.
Ordinary zeta series and Euler product in their stated
\(\Rea s>1\) domain; consulted 12 November 2026.

\bibitem{dlmfpnt}
NIST Digital Library of Mathematical Functions,
\href{https://dlmf.nist.gov/27.2#E3}{Equation 27.2.3, prime number theorem}.
Consulted 12 November 2026.

\bibitem{collision}
Edward Baker project, with acknowledged LLM assistance,
\emph{Newman collision reductions}, living internal manuscript.
\repo{papers/quasi-rh/quasi-rh-exponent-descent/newman_collisions/newman_collision_reductions.tex}{Canonical TeX source}.
{\raggedright Relevant labels: \texttt{prop:finitecollision}, \texttt{cor:rhcollision},
\texttt{prop:absmass}, \texttt{sc:derivativemass},
\texttt{sc:character}, \texttt{sc:remaining}, \texttt{lj:mirror},
\texttt{lj:paidtest}, and \texttt{lj:quartic}.\par}

\bibitem{hda}
Edward Baker project, with acknowledged LLM assistance,
\emph{Complex-strip remainder and buffered crossover}, HDA Ansatz
Note04, 25 October 2026, especially CS3 and CS6.
\repo{papers/investigations-5/01_heat_density_asymptotics/ansatz/notes/04_COMPLEX_STRIP_REMAINDER_AND_BUFFERED_CROSSOVER_20261025.md}{Proof note};
\repo{papers/investigations-5/01_heat_density_asymptotics/ansatz/reviews/03_COMPLEX_STRIP_AND_BUFFERED_CROSSOVER_ASSESSMENT_20261025.md}{same-assistant assessment}.

\bibitem{gnpe01}
Edward Baker project, with acknowledged LLM assistance,
\emph{GNP-RS-E01 r1: uniform outer Green baseline},
11 November 2026, especially E2 and E9--E19.
\repo{papers/investigations-5/09_green_potentials_newman_flow/01_riemann_siegel_evolution/notes/03_GNP_RS_E01_UNIFORM_OUTER_BASELINE_PROOF_20261111.md}{Proof and approximation ledger}.

\bibitem{s04}
Edward Baker project, with acknowledged LLM assistance,
\emph{Phase-retaining full-cutoff scaling}, S04 Note01,
18 October 2026, and \emph{Uniform local residual transfer and cutoff
payment obstruction}, S04 Note03, 19 October 2026.
\repo{papers/investigations-3/11_arithmetic_collision_rigidity/symmetries/04_time_height_scaling/notes/01_PHASE_RETAINING_FULL_CUTOFF_SCALING_20261018.md}{Note01};
\repo{papers/investigations-3/11_arithmetic_collision_rigidity/symmetries/04_time_height_scaling/notes/03_UNIFORM_LOCAL_RESIDUAL_TRANSFER_AND_CUTOFF_PAYMENT_OBSTRUCTION_20261019.md}{Note03}.

\bibitem{r1}
Drafted for Edward Baker with substantial OpenAI Codex assistance,
\emph{RSCT-B01: a reflected Poisson constraint and its unpaid arithmetic sign},
11 November 2026.
\rsctnote{03_RSCT_B01_REFLECTED_POISSON_TEST_20261111.md}{RSCT Note03 (R1)}.

\bibitem{r2}
Same project and assistance,
\emph{Signed three-height estimate}, 11 November 2026.
\rsctnote{04_RSCT_B01_R2_SIGNED_THREE_HEIGHT_ESTIMATE_20261111.md}{RSCT Note04 (R2)}.

\bibitem{r3}
Same project and assistance,
\emph{Event-subtracted Pick test}, 11 November 2026.
\rsctnote{05_RSCT_B01_R3_EVENT_SUBTRACTED_PICK_TEST_20261111.md}{RSCT Note05 (R3)}.

\bibitem{r4}
Same project and assistance,
\emph{Microscopic event expansion and the complete population cost},
11 November 2026.
\rsctnote{06_RSCT_B01_R4_MICROSCOPIC_EVENT_EXPANSION_20261111.md}{RSCT Note06 (R4)}.

\bibitem{r5}
Same project and assistance,
\emph{Signed correlation disk test}, 11 November 2026.
\rsctnote{07_RSCT_B01_R5_SIGNED_CORRELATION_DISK_TEST_20261111.md}{RSCT Note07 (R5)}.

\bibitem{r6}
Same project and assistance,
\emph{Exact circle test}, 11 November 2026.
\rsctnote{08_RSCT_B01_R6_EXACT_CIRCLE_TEST_20261111.md}{RSCT Note08 (R6)}.

\bibitem{r7}
Same project and assistance,
\emph{Multiplicative phase test}, 11 November 2026.
\rsctnote{09_RSCT_B01_R7_MULTIPLICATIVE_PHASE_TEST_20261111.md}{RSCT Note09 (R7)}.

\bibitem{r8}
Same project and assistance,
\emph{Growing prime calibration}, 11 November 2026.
\rsctnote{10_RSCT_B01_R8_GROWING_PRIME_CALIBRATION_20261111.md}{RSCT Note10 (R8)}.

\bibitem{r9}
Same project and assistance,
\emph{Mirror-jet compatibility}, 11 November 2026.
\rsctnote{11_RSCT_B01_R9_MIRROR_JET_COMPATIBILITY_20261111.md}{RSCT Note11 (R9)}.

\bibitem{gnp19}
Edward Baker project, with acknowledged LLM assistance,
\emph{Manuscript and averaged Green values}, GNP-RS Note19,
11 November 2026.
\repo{papers/investigations-5/09_green_potentials_newman_flow/01_riemann_siegel_evolution/notes/19_GNP_RS_MANUSCRIPT_AND_AVERAGED_GREEN_VALUES_20261111.md}{Proof, scope and same-assistant evidence}.

\bibitem{gnp20}
Edward Baker project, with acknowledged LLM assistance,
\emph{Green geometry and critical time}, GNP-RS Note20,
11 November 2026.
\repo{papers/investigations-5/09_green_potentials_newman_flow/01_riemann_siegel_evolution/notes/20_GNP_RS_GREEN_GEOMETRY_AND_CRITICAL_TIME_20261111.md}{Geometric distinction and controls}.
\end{thebibliography}

\end{document}
